\documentclass[11pt]{article}

\usepackage[margin=1in]{geometry}
\usepackage{amsmath,amssymb,bm,amsthm}
\usepackage{booktabs}
\usepackage{tikz}
\usetikzlibrary{arrows.meta,positioning,calc}
\usepackage{microtype}
\usepackage[hidelinks]{hyperref}

\newtheorem{proposition}{Proposition}
\theoremstyle{remark}
\newtheorem{remark}{Remark}

\title{Working Notes: Plug-and-Play Twin-Field / Phase-Matching QKD}
\author{}
\date{Updated 7 October 2026}

\begin{document}

\maketitle

\section{Current picture}

The project has converged to a fairly clear security story.

The original guessing-probability / min-entropy route appears too loose for the
protocol and has been set aside. The current route is a phase-error SDP based
on the Gram-matrix method of arXiv:1901.01942.  Alice and Bob are allowed to
receive an arbitrary optical state supplied by the untrusted central node and
entangled with Eve; the trusted structure is local phase modulation together
with whatever attenuation and monitoring assumptions are stated explicitly.

Stages~1--3 establish a negative result for mean-only source information.
A bound on average photon number, even combined with very strong trusted
attenuation, leaves a concentration degree of freedom: Eve may use mostly
vacuum or weak pulses together with rare bright pulses.  As the central pass
probability decreases with distance, the energy of the rare component
compatible with the same first-moment bound can grow roughly as
\(1/P_{\rm pass}\).  Strong attenuation changes the overlap kernel and makes
the problem numerically tractable through the rescaled variable
\(u=\tau m\), but it does not remove this concentration mechanism.

Stage~4 isolates the source information that does remove that freedom.  It is
an \emph{abstract, hardware-independent} conditional upper-tail statement:
among monitor-accepted rounds, all but an \(\epsilon_{\rm mon}\) fraction have
effective delivered energies below fixed caps \(u_{A,+},u_{B,+}\).  The
current Stage-4 SDP is rigorous conditional on that statement and gives
positive conditional key rates through the largest distance presently scanned
(\(40\) km for the illustrative parameters).  It does \emph{not} yet derive
the monitor statement from a physical photon-number analyzer.

The next step is therefore a physical monitor model rather than a new
phase-error formalism.  Stage~5A will treat a trusted beam splitter and
photon-number analyzer jointly, allowing the same beam splitter both to send a
large fraction of the incoming light to the monitor and to leave a weak signal
for encoding.  Stage~5B will add optional trusted attenuation after that
joint monitor.  The extra attenuator is an implementation/generalisation
choice, not something required to create the common phase orbit.

\section{Protocol and security model}

The central node is untrusted and supplies the optical pulses. Alice and Bob
apply phase modulation and send the light back to the node for interference
and detection. Eve is therefore allowed to control both the source and the
measurement device; Alice and Bob trust only their local modulation,
attenuation, and monitoring.

The phase encoding uses \(M\) bases,
\[
    \phi_x = \frac{x\pi}{M},
    \qquad x\in\{0,\ldots,M-1\},
\]
with the bit value contributing an additional phase shift \(a\pi\). In the
present phase-error formulation, basis \(x=0\) is used for key generation and
the remaining \(M-1\) bases are test bases.

The working key-rate expression is
\begin{equation}
    R = P_{\rm pass}
    \left[
        1-h_2(e_{\rm ph})
        -f_{\rm ec}h_2(e_{\rm bit})
    \right].
\end{equation}
For the current numerical choice
\[
    e_{\rm bit}=0.015,
    \qquad
    f_{\rm ec}=1.15,
\]
positive key requires approximately
\[
    e_{\rm ph}<0.292.
\]

The phase-error objective is the analogue of Eq.~(A7) of
arXiv:1901.01942: \(e_{\rm ph}\) is an affine function of off-diagonal
entries of the Gram matrix of Eve's post-measurement states. The test bases
are essential. Without them the complementary-basis phase error is
unconstrained and the SDP returns the trivial value \(e_{\rm ph}=1/2\).

\section{Derivation of the phase-error objective for an arbitrary untrusted source}
\label{sec:a7-general}

Equation~(A7) of arXiv:1901.01942 was written there for coherent-state
phase encoding.  For the present protocol, however, the same algebraic
formula does not require the key states to be coherent, product, or even
prepared by Alice and Bob.  It follows from source replacement, the
announcement-dependent bit correction, and the definition of phase error on
the virtual qubits.

Let
\[
    P_A=e^{i\pi\hat n_A},
    \qquad
    P_B=e^{i\pi\hat n_B}.
\]
Eve may prepare an arbitrary pure state
\[
    |\Psi\rangle_{ABE_0},
\]
where \(E_0\) is retained by Eve.  The four key-basis states generated by
Alice's and Bob's bits \(a,b\in\{0,1\}\) are
\begin{equation}
    |\psi_{ab}\rangle
    =
    (P_A^a\otimes P_B^b\otimes I_{E_0})
    |\Psi\rangle.
    \label{eq:a7-key-states-general}
\end{equation}
No assumption on the photon-number distribution or on factorisation between
\(A\), \(B\), and \(E_0\) is made.

Define the \(X\)-basis labels
\[
    |x_0\rangle=|+\rangle,
    \qquad
    |x_1\rangle=|-\rangle.
\]
A source-replacement state for the actual random bit choices is
\begin{equation}
    |\Omega\rangle_{A_0B_0ABE_0}
    =
    \frac12
    \sum_{a,b=0}^1
    |x_a x_b\rangle_{A_0B_0}
    |\psi_{ab}\rangle_{ABE_0}.
    \label{eq:a7-source-replacement}
\end{equation}
Indeed,
\begin{equation}
    \operatorname{Tr}_{A_0B_0}
    |\Omega\rangle\!\langle\Omega|
    =
    \frac14
    \sum_{a,b}
    |\psi_{ab}\rangle\!\langle\psi_{ab}|,
    \label{eq:a7-eve-equivalence}
\end{equation}
which is exactly the state seen by Eve in the prepare-and-measure protocol
when the key bits are independent and uniform.  Measuring \(A_0B_0\) in the
\(X\) basis can therefore be delayed until after Eve's announcement.

Let Eve's most general attack, including the untrusted node, be represented
by an isometry
\begin{equation}
    V|\psi_{ab}\rangle
    =
    |e^+_{ab}\rangle_E|+\rangle_Z
    +
    |e^-_{ab}\rangle_E|-\rangle_Z
    +
    |e^\varnothing_{ab}\rangle_E|\varnothing\rangle_Z,
    \label{eq:a7-attack-isometry}
\end{equation}
where the vectors on \(E\) are sub-normalised and the orthogonal register
\(Z\) stores the public announcement.  The average key-basis pass probability
is
\begin{equation}
    P_{\rm pass}
    =
    \frac14
    \sum_{a,b}
    \left(
      \|e^+_{ab}\|^2+\|e^-_{ab}\|^2
    \right).
    \label{eq:a7-ppass-general}
\end{equation}

After projecting onto the successful outcomes, the normalised virtual state
is
\begin{equation}
    |\Omega_{\rm pass}\rangle
    =
    \frac{1}{2\sqrt{P_{\rm pass}}}
    \sum_{a,b}
    |x_a x_b\rangle
    \left(
      |e^+_{ab}\rangle|+\rangle_Z
      +
      |e^-_{ab}\rangle|-\rangle_Z
    \right).
    \label{eq:a7-pass-state}
\end{equation}
The protocol uses \(\Psi^+\) as the target Bell state.  For a \(+\)
announcement Bob does nothing; for a \(-\) announcement he flips his raw bit.
On the virtual qubit this is the Pauli \(Z_{B_0}\), since
\[
    Z|+\rangle=|-\rangle,
    \qquad
    Z|-\rangle=|+\rangle.
\]
After this controlled correction,
\begin{equation}
\begin{aligned}
    |\widetilde\Omega_{\rm pass}\rangle
    =
    \frac{1}{2\sqrt{P_{\rm pass}}}
    \Big[&
    |++\rangle |F_{++}\rangle
    +|+-\rangle |F_{+-}\rangle\\
    &+|-+\rangle |F_{-+}\rangle
    +|--\rangle |F_{--}\rangle
    \Big],
    \label{eq:a7-corrected-state}
\end{aligned}
\end{equation}
where
\begin{align}
    |F_{++}\rangle
    &=
    |e^+_{++}\rangle|+\rangle_Z
    +
    |e^-_{+-}\rangle|-\rangle_Z,\\
    |F_{+-}\rangle
    &=
    |e^+_{+-}\rangle|+\rangle_Z
    +
    |e^-_{++}\rangle|-\rangle_Z,\\
    |F_{-+}\rangle
    &=
    |e^+_{-+}\rangle|+\rangle_Z
    +
    |e^-_{--}\rangle|-\rangle_Z,\\
    |F_{--}\rangle
    &=
    |e^+_{--}\rangle|+\rangle_Z
    +
    |e^-_{-+}\rangle|-\rangle_Z.
    \label{eq:a7-F-vectors}
\end{align}

The target state \(|\Psi^+\rangle\) is a \(+1\) eigenstate of
\(Y\otimes Y\).  Hence, if the phase error is the probability of obtaining
the wrong \(Y\)-basis parity,
\begin{equation}
    e_{\rm ph}
    =
    \frac{1-\langle Y_{A_0}\otimes Y_{B_0}\rangle}{2}.
    \label{eq:a7-phase-definition}
\end{equation}
In the \(X\) basis,
\[
    Y|+\rangle=-i|-\rangle,
    \qquad
    Y|-\rangle=i|+\rangle,
\]
and therefore
\begin{align}
    Y\otimes Y|++\rangle&=-|--\rangle,\\
    Y\otimes Y|--\rangle&=-|++\rangle,\\
    Y\otimes Y|+-\rangle&=|-+\rangle,\\
    Y\otimes Y|-+\rangle&=|+-\rangle.
\end{align}
Substituting Eq.~\eqref{eq:a7-corrected-state} gives
\begin{equation}
\begin{aligned}
    \langle Y\otimes Y\rangle
    =
    \frac{1}{2P_{\rm pass}}
    \operatorname{Re}\Big[
       \langle F_{+-}|F_{-+}\rangle
       -
       \langle F_{++}|F_{--}\rangle
    \Big].
    \label{eq:a7-yy-F}
\end{aligned}
\end{equation}
Because the announcement states \(|+\rangle_Z\) and \(|-\rangle_Z\) are
orthogonal,
\begin{align}
    \langle F_{++}|F_{--}\rangle
    &=
    \langle e^+_{++}|e^+_{--}\rangle
    +
    \langle e^-_{+-}|e^-_{-+}\rangle,\\
    \langle F_{+-}|F_{-+}\rangle
    &=
    \langle e^+_{+-}|e^+_{-+}\rangle
    +
    \langle e^-_{++}|e^-_{--}\rangle.
\end{align}
Combining these expressions with
Eq.~\eqref{eq:a7-phase-definition} yields
\begin{equation}
\boxed{
\begin{aligned}
    e_{\rm ph}
    =
    \frac12
    +
    \frac{1}{4P_{\rm pass}}
    \operatorname{Re}\Big[
      &\langle e^+_{++}|e^+_{--}\rangle
      -
      \langle e^-_{++}|e^-_{--}\rangle\\
      &-
      \langle e^+_{+-}|e^+_{-+}\rangle
      +
      \langle e^-_{+-}|e^-_{-+}\rangle
    \Big].
\end{aligned}}
    \label{eq:a7-general}
\end{equation}
This is exactly Eq.~(A7) of arXiv:1901.01942.

\begin{proposition}[Generality of Eq.~(A7)]
Equation~\eqref{eq:a7-general} does not use coherent-state amplitudes,
Poisson photon statistics, product states, or trusted sources.  It is valid
for any four pure key-basis states \(\{|\psi_{ab}\rangle\}\), including the
joint states generated by local phase modulation of an arbitrary source
entangled with Eve, provided that the bits are uniform, the public successful
announcements are \(+\) and \(-\), and Bob applies the stated correction.
\end{proposition}

For the attenuated stages the states returned to Eve are mixed after the
trusted loss ports are traced out.  Choose the common purification introduced
in Section~\ref{sec:stage2-rigorous} and include its purifying register in the
residual vectors \(|e^z_{ab}\rangle\).  The derivation above is unchanged.
Optimising over arbitrary Gram matrices on that enlarged purifying space can
only enlarge the set of attacks relative to the physical one, so the resulting
phase-error optimum remains conservative.

Thus the coherent-state restriction of the original Appendix~A is not a
restriction on the algebraic phase-error objective used here.  The source
model affects the SDP through the admissible overlaps and observed
probabilities, not through the form of Eq.~\eqref{eq:a7-general}.


\section{Route that was abandoned: guessing probability}

The guessing-probability SDP was corrected in two important ways.

First, the photon-number tail / energy constraint was fixed so that probability
mass outside the numerical cutoff is charged consistently to the combined
energy budget. Second, Eve's final bit guess was allowed to depend on the
basis after that basis is publicly announced, while her earlier announcement
remains basis-independent.

After those corrections, for
\[
    \bar n = 10^{-2},
    \qquad L=0,
    \qquad M=4,
\]
the SDP gives
\[
    P_{\rm guess}=0.9336.
\]
A positive key from that min-entropy bound would require
\[
    P_{\rm guess}<0.9145.
\]
Thus the corrected guessing-probability route gives zero key already at zero
distance. Since the phase-error calculation gives positive key in the same
regime, this is interpreted as looseness of the min-entropy bound rather than
as evidence that the protocol itself is insecure.

\section{Validation of the phase-error implementation}

Two checks have been used.

\begin{enumerate}
    \item The implementation reproduces the trusted-source finite-test-state
    phase-matching calculation of Fig.~6(a) of arXiv:1901.01942. The
    reproduced PLOB crossing is about \(48\) dB, compared with roughly
    \(50\) dB in the published figure.

    \item If the plug-and-play source is artificially pinned to the same
    coherent-state distribution, two independent implementations agree to
    about five digits. This checks the overlap construction, the loss
    convention, the truncation, and the indexing in the phase-error objective.
\end{enumerate}

\section{Stage 1: untrusted source with only mean-energy monitoring}
\label{sec:stage1-rigorous}

This section writes the bare-source model as an exact infinite-dimensional
optimization first, and only afterwards introduces the finite-dimensional
relaxation used numerically. Keeping those two levels separate is important:
the SDP is rigorous only if every physical source and attack is mapped to a
feasible point of the numerical problem.

\subsection{Exact source overlaps}

Let \(A\) and \(B\) denote the optical modes entering Alice's and Bob's
trusted labs and let \(E_0\) purify the state supplied by Eve. Thus, without
loss of generality, the source can be written as a pure state
\[
    \lvert\Psi\rangle_{ABE_0},
    \qquad
    \rho_{AB}=\operatorname{Tr}_{E_0}\lvert\Psi\rangle\!\langle\Psi\rvert.
\]
For a signal label \(s\), Alice and Bob apply
\[
    U_s
    =
    \exp\!\left[
      i\phi_A^s \hat n_A+i\phi_B^s\hat n_B
    \right],
\]
and the global encoded state is
\[
    \lvert\psi_s\rangle
    =
    (U_s\otimes I_{E_0})\lvert\Psi\rangle.
\]

\begin{proposition}[Exact overlap formula]
For any two signal labels \(s,t\),
\begin{equation}
    \Lambda_{st}
    :=
    \langle\psi_s\vert\psi_t\rangle
    =
    \sum_{n_A,n_B\ge0}
      p_{n_A,n_B}
      e^{i\Delta\phi_A n_A+i\Delta\phi_B n_B},
    \label{eq:stage1-exact-overlap}
\end{equation}
where
\[
    p_{n_A,n_B}
    :=
    \langle n_A,n_B\vert\rho_{AB}\vert n_A,n_B\rangle,
\]
\(\Delta\phi_X=\phi_X^t-\phi_X^s\), and
\(\sum_{n_A,n_B}p_{n_A,n_B}=1\).
Equation~\eqref{eq:stage1-exact-overlap} is exact even if
\(\rho_{AB}\) has arbitrary Fock-basis coherences.
\end{proposition}

\begin{proof}
Since \(U_s^\dagger U_t\) is diagonal in the joint Fock basis,
\begin{align}
    \Lambda_{st}
    &=
    \operatorname{Tr}
    \left[
      \rho_{AB}U_s^\dagger U_t
    \right] \\
    &=
    \sum_{n_A,n_B}
    \langle n_A,n_B\vert\rho_{AB}\vert n_A,n_B\rangle
    e^{i\Delta\phi_A n_A+i\Delta\phi_B n_B}.
\end{align}
All off-diagonal matrix elements of \(\rho_{AB}\) vanish from the trace.
\end{proof}

Thus the unknown source enters the Gram constraints only through a classical
probability distribution \(p_{n_A,n_B}\), even though the physical source
itself need not be phase randomized.

\subsection{Exact Gram constraints}

Represent Eve's channel and announcement by an isometry
\[
    V:\lvert\psi_s\rangle
    \longmapsto
    \sum_z \lvert z\rangle_Z\lvert e_s^z\rangle_E,
\]
where \(z\) is the public announcement and the
\(\lvert e_s^z\rangle\) are sub-normalized residual states. Isometry
preserves inner products, so
\begin{equation}
    \Lambda_{st}
    =
    \sum_z
      \langle e_s^z\vert e_t^z\rangle.
    \label{eq:stage1-gram-conservation}
\end{equation}
If \(G\) is the Gram matrix of all vectors
\(\{\lvert e_s^z\rangle\}_{s,z}\), then \(G\succeq0\),
Eq.~\eqref{eq:stage1-gram-conservation} is linear in \(G\), and the observed
probabilities are diagonal entries,
\[
    P(z\mid s)
    =
    \lVert e_s^z\rVert^2
    =
    G[(z,s),(z,s)].
\]
Consequently the pass probabilities, QBER constraints, and the phase-error
objective inherited from Eq.~(A7) of arXiv:1901.01942 are affine functions of
\(G\).

The exact optimization would therefore maximize the phase-error functional
over \(G\) and the infinite distribution \(p\), subject to
\[
    G\succeq0,
    \qquad
    p_{n_A,n_B}\ge0,
    \qquad
    \sum p_{n_A,n_B}=1,
\]
the overlap equalities
\eqref{eq:stage1-exact-overlap}--\eqref{eq:stage1-gram-conservation}, the
observed statistics, and the monitored energy constraints
\begin{equation}
    \mathbb E[n_A]\le\frac{\bar n}{2},
    \qquad
    \mathbb E[n_B]\le\frac{\bar n}{2}.
    \label{eq:stage1-energy-exact}
\end{equation}

\subsection{Rigorous photon-number truncation}

For computation define
\[
    S_N
    :=
    \{(n_A,n_B):n_A+n_B\le N\}
\]
and
\[
    \varepsilon_N
    :=
    1-\sum_{(n_A,n_B)\in S_N}p_{n_A,n_B}.
\]
Write
\[
    \Lambda_{st}
    =
    \Lambda^{(N)}_{st}+r_{st},
\]
where
\[
    \Lambda^{(N)}_{st}
    =
    \sum_{(n_A,n_B)\in S_N}
    p_{n_A,n_B}
    e^{i\Delta\phi_A n_A+i\Delta\phi_B n_B}.
\]
Because every discarded phase factor has modulus one,
\begin{equation}
    \lvert r_{st}\rvert
    \le
    \varepsilon_N.
    \label{eq:stage1-tail-overlap}
\end{equation}
Hence every physical attack satisfies
\begin{equation}
    \left|
      \sum_zG[(z,s),(z,t)]
      -
      \Lambda^{(N)}_{st}
    \right|
    \le
    \varepsilon_N.
    \label{eq:stage1-soc}
\end{equation}

The tail must also be charged to the energy budget. Since
\(n_A+n_B\ge N+1\) on \(S_N^c\),
\begin{equation}
    \sum_{S_N}p_{n_A,n_B}(n_A+n_B)
    +(N+1)\varepsilon_N
    \le
    \bar n.
    \label{eq:stage1-combined-energy}
\end{equation}
Dropping the nonnegative tail also gives the individually valid constraints
\begin{equation}
    \sum_{S_N}p_{n_A,n_B}n_A
    \le\frac{\bar n}{2},
    \qquad
    \sum_{S_N}p_{n_A,n_B}n_B
    \le\frac{\bar n}{2}.
    \label{eq:stage1-perarm-energy}
\end{equation}

Equations~\eqref{eq:stage1-tail-overlap}--\eqref{eq:stage1-perarm-energy}
define an \emph{outer relaxation}: every physical source obeying
\eqref{eq:stage1-energy-exact} maps to a feasible truncated point. Therefore
maximizing \(e_{\rm ph}\) over this relaxation gives a valid upper bound.  The
phase-error functional itself is justified for the arbitrary source model in
Section~\ref{sec:a7-general}.

\begin{remark}[Why the tail cannot be charged separately to both arms]
It would be unsound to add \((N+1)\varepsilon_N\) to each of the two
per-arm constraints. A discarded point such as \((0,N+1)\) has
\(n_A=0\), so there is no lower bound \(n_A\ge N+1\) on the discarded set.
Only the combined photon number has the required lower bound.
\end{remark}

\subsection{A first moment cannot control bunching}

For a one-mode distribution with fixed mean \(\mu>0\), consider
\begin{equation}
    P(N=0)=1-\frac{\mu}{K},
    \qquad
    P(N=K)=\frac{\mu}{K},
    \qquad K\ge\mu.
    \label{eq:stage1-two-point}
\end{equation}
Then \(\mathbb E[N]=\mu\), but
\begin{equation}
    g^{(2)}
    =
    \frac{\mathbb E[N(N-1)]}{\mathbb E[N]^2}
    =
    \frac{K-1}{\mu},
    \label{eq:stage1-unbounded-g2}
\end{equation}
which diverges as \(K\to\infty\). Thus no bound on \(g^{(2)}\), or on any
higher moment, follows from a mean-energy constraint alone.

This does not by itself prove the optimal QKD attack, but it proves the
structural point used later: the monitored first moment leaves arbitrarily
strong pulse-to-pulse concentration available to Eve.

Numerically, at
\[
    \bar n=10^{-2},
    \qquad L=0,
    \qquad N_{\max}=7,
\]
the current phase-error values are
\begin{center}
\begin{tabular}{ccc}
\toprule
\(M\) & number of test states & \(e_{\rm ph}\) \\
\midrule
2 & 2 & 0.5000 \\
3 & 4 & 0.2865 \\
4 & 6 & 0.2573 \\
5 & 8 & 0.2572 \\
\bottomrule
\end{tabular}
\end{center}
so the test bases constrain the phase structure, but the distance remains
limited by source statistics. A representative optimizer has
\(g^{(2)}\sim36\).

\section{Stage 2: trusted attenuation as a rigorous Gram-matrix relaxation}
\label{sec:stage2-rigorous}

Let \(\mathcal L_\tau\) denote a trusted pure-loss channel of transmissivity
\[
    \tau=t^2.
\]
For simplicity the same \(\tau\) is written for both arms; all formulas extend
immediately to \(\tau_A\neq\tau_B\).

The purpose of this section is to make precise what trusted attenuation buys.
There are two logically distinct gains:

\begin{enumerate}
    \item \emph{Structural gain.}  Because pure loss is phase covariant, all
    post-loss encoded states lie on a common unitary orbit.  This lets us
    construct one consistent pure-state family and hence recover a Gram-matrix
    description even though the physical states returned to Eve are mixed.

    \item \emph{Constraint gain.}  The post-loss photon statistics are not free:
    they must lie in the image of binomial thinning.  Equivalently, the overlap
    kernel changes from a pure phase factor to
    \((1-\tau+\tau e^{i\Delta\phi})^m\).  This is the actual source restriction
    introduced by the trusted attenuator.
\end{enumerate}

The first point restores tractability; the second is what can tighten the
security bound.

\subsection{Physical model before and after the attenuator}

Let Eve prepare an arbitrary pure state
\[
    |\Psi\rangle_{ABE_0},
\]
where \(A,B\) are the optical systems entering Alice's and Bob's laboratories
and \(E_0\) is retained by Eve.  Alice and Bob encode signal label \(s\) by
\[
    U_s
    =
    e^{i\phi_A^s\hat n_A+i\phi_B^s\hat n_B}.
\]
The actual physical order is
\[
    |\Psi\rangle
    \xrightarrow{U_s}
    (U_s\otimes I_{E_0})|\Psi\rangle
    \xrightarrow{\mathcal L_\tau^A\otimes\mathcal L_\tau^B}
    \rho_s^{TE_0},
\]
where \(T=(T_A,T_B)\) denotes the transmitted optical modes and the trusted
reflected ports have been discarded.

Pure loss is phase covariant:
\begin{equation}
    \mathcal L_\tau
    \!\left(
      e^{i\theta\hat n}\rho e^{-i\theta\hat n}
    \right)
    =
    e^{i\theta\hat n}
    \mathcal L_\tau(\rho)
    e^{-i\theta\hat n}.
    \label{eq:stage2-covariance}
\end{equation}
Hence we may equivalently regard loss as occurring before the phase rotation.
Define
\begin{equation}
    \rho_0^{TE_0}
    :=
    \left(
      \mathcal L_\tau^A\otimes\mathcal L_\tau^B\otimes I_{E_0}
    \right)
    \left(
      |\Psi\rangle\!\langle\Psi|
    \right).
    \label{eq:stage2-base-mixed}
\end{equation}
Then every encoded state returned to the untrusted node satisfies
\begin{equation}
    \rho_s^{TE_0}
    =
    W_s\rho_0^{TE_0}W_s^\dagger,
    \qquad
    W_s
    =
    e^{i\phi_A^s\hat n_{T_A}+i\phi_B^s\hat n_{T_B}}
    \otimes I_{E_0}.
    \label{eq:stage2-unitary-orbit}
\end{equation}

Thus the family is not an arbitrary collection of mixed states: it is one
mixed state acted on by a known representation of the phase group.

\subsection{What the common orbit buys: a consistent purification}

Choose one purification
\[
    |\chi_0\rangle_{TE_0R}
\]
of \(\rho_0^{TE_0}\), where \(R\) is a fictitious reference, and define
\begin{equation}
    |\chi_s\rangle
    :=
    (W_s\otimes I_R)|\chi_0\rangle.
    \label{eq:stage2-common-purification}
\end{equation}
Then \(|\chi_s\rangle\) purifies \(\rho_s^{TE_0}\) for every \(s\).

\begin{proposition}[Purified overlap after trusted loss]
The purified family obeys
\begin{align}
    \langle\chi_s|\chi_t\rangle
    &=
    \operatorname{Tr}_{TE_0}
    \left[
      \rho_0^{TE_0}W_s^\dagger W_t
    \right]
    \nonumber\\
    &=
    \operatorname{Tr}_{T}
    \left[
      \rho_{0,T}\,
      e^{i\Delta\phi_A\hat n_{T_A}
        +i\Delta\phi_B\hat n_{T_B}}
    \right],
    \label{eq:stage2-purified-overlap}
\end{align}
where
\[
    \rho_{0,T}
    :=
    \operatorname{Tr}_{E_0}\rho_0^{TE_0}.
\]
\end{proposition}

\begin{proof}
Use the purification identity
\[
    \langle\chi_0|(X\otimes I_R)|\chi_0\rangle
    =
    \operatorname{Tr}(\rho_0X).
\]
Since \(W_s^\dagger W_t\) acts trivially on \(E_0\), the partial trace over
\(E_0\) gives the second line.
\end{proof}

The notation \(\rho_{0,T}\) is deliberately used instead of a superscript
\(T\): no matrix transpose is involved.

This is the precise point at which our setting differs from the generic lossy
channel-discrimination setting of arXiv:2012.11104.  There, a pure input can
become a family of mixed states with no common known eigenbasis, so the
original pure-state Gram construction no longer applies directly.  Here,
phase covariance supplies the extra unitary-orbit structure needed to choose a
single compatible purification for all signal labels.

\subsection{From the purified family to the QKD Gram matrix}

Let the untrusted node implement an arbitrary physical instrument on \(TE_0\)
with Kraus operators \(\{K_z\}_z\), where \(z\) labels the public
announcement and
\[
    \sum_zK_z^\dagger K_z=I_{TE_0}.
\]
On the purified system define the sub-normalized residual vectors
\begin{equation}
    |e_s^z\rangle
    :=
    (K_z\otimes I_R)|\chi_s\rangle.
    \label{eq:stage2-residual-vectors}
\end{equation}
Then
\[
    \|e_s^z\|^2=P(z|s),
\]
and
\begin{align}
    \sum_z\langle e_s^z|e_t^z\rangle
    &=
    \sum_z
    \langle\chi_s|
    (K_z^\dagger K_z\otimes I_R)
    |\chi_t\rangle
    \nonumber\\
    &=
    \langle\chi_s|\chi_t\rangle.
    \label{eq:stage2-gram-conservation}
\end{align}

Therefore the same Gram-matrix structure used in Stage~1 survives:
\begin{equation}
    \sum_z G[(z,s),(z,t)]
    =
    \Lambda^{(\tau)}_{st},
    \label{eq:stage2-gram-overlap}
\end{equation}
where
\[
    \Lambda^{(\tau)}_{st}
    :=
    \langle\chi_s|\chi_t\rangle.
\]

There is one deliberate relaxation.  In the physical model, all residual
vectors must arise from one common instrument \(K_z\otimes I_R\), with the
fictitious reference \(R\) untouched.  In the SDP we enforce only positivity
of the Gram matrix together with
Eq.~\eqref{eq:stage2-gram-overlap} and the observed diagonal entries.  This
can admit Gram matrices that do not arise from any physical instrument acting
trivially on \(R\).  Hence
\[
    \mathcal F_{\rm physical}
    \subseteq
    \mathcal F_{\rm Gram},
\]
and, since the SDP maximizes \(e_{\rm ph}\),
\begin{equation}
    e_{\rm ph}^{\rm physical}
    \le
    e_{\rm ph}^{\rm Gram}.
    \label{eq:stage2-outer-direction}
\end{equation}
Thus the common-purification construction is used as a conservative
\emph{outer relaxation}; it need not reproduce Eve's exact physical
information.

\begin{remark}[Why keeping the physical reflected port is not useful here]
If the trusted reflected loss ports are retained, the beamsplitter dilation of
loss is unitary and independent of \(s\).  The global overlaps are then just
the pre-loss overlaps, so attenuation produces no new overlap constraint.
The useful construction is to trace out the trusted loss ports first, obtaining
the physical mixed transmitted state, and only then introduce a mathematical
purification of that reduced state.
\end{remark}

\subsection{Exact binomial-thinning kernel}

Let
\[
    p_{m_A,m_B}
    =
    \langle m_A,m_B|\rho_{AB}|m_A,m_B\rangle
\]
be the incoming joint photon-number distribution before trusted attenuation.
For one input Fock component \(m\),
\[
    N|m\sim\operatorname{Binomial}(m,\tau).
\]
Consequently
\begin{align}
    \mathbb E[e^{i\Delta\phi N}|m]
    &=
    \sum_{n=0}^{m}
      \binom mn
      \tau^n(1-\tau)^{m-n}
      e^{i\Delta\phi n}
    \nonumber\\
    &=
    \left(
      1-\tau+\tau e^{i\Delta\phi}
    \right)^m.
    \label{eq:stage2-binomial-kernel}
\end{align}
Define
\[
    z_X^{st}
    :=
    1-\tau+\tau e^{i\Delta\phi_X^{st}}
    =
    1-\tau(1-e^{i\Delta\phi_X^{st}}).
\]
Then the exact purified overlap is
\begin{equation}
    \boxed{
    \Lambda^{(\tau)}_{st}
    =
    \sum_{m_A,m_B\ge0}
      p_{m_A,m_B}
      (z_A^{st})^{m_A}
      (z_B^{st})^{m_B}.
    }
    \label{eq:stage2-exact-kernel}
\end{equation}

Equivalently, if \(q_{n_A,n_B}\) denotes the transmitted photon-number
distribution, then
\begin{equation}
    q_{n_A,n_B}
    =
    \sum_{\substack{m_A\ge n_A\\m_B\ge n_B}}
      p_{m_A,m_B}
      B_\tau(n_A|m_A)
      B_\tau(n_B|m_B),
    \label{eq:stage2-thinning-map}
\end{equation}
where
\[
    B_\tau(n|m)
    =
    \binom mn\tau^n(1-\tau)^{m-n},
\]
and
\[
    \Lambda^{(\tau)}_{st}
    =
    \sum_{n_A,n_B}
      q_{n_A,n_B}
      e^{i\Delta\phi_A^{st}n_A+i\Delta\phi_B^{st}n_B}.
\]
Equation~\eqref{eq:stage2-exact-kernel} simply eliminates \(q\) analytically.

This makes clear what trusted attenuation constrains:
\[
    q=B_\tau p.
\]
The transmitted distribution is not an arbitrary distribution with the right
mean; it must be obtainable by binomial thinning of some admissible incoming
distribution.

\subsection{Exact infinite-dimensional Stage-2 optimization}

Before any photon-number cutoff, the trusted-attenuation problem can be
written as
\begin{equation}
\begin{aligned}
    \text{maximize}_{G,p}\quad&
    e_{\rm ph}(G)
    \\
    \text{subject to}\quad&
    G\succeq0,
    \\
    &p_{m_A,m_B}\ge0,
    \qquad
    \sum_{m_A,m_B}p_{m_A,m_B}=1,
    \\
    &\sum_zG[(z,s),(z,t)]
    =
    \sum_{m_A,m_B}
      p_{m_A,m_B}
      (z_A^{st})^{m_A}(z_B^{st})^{m_B}
    \quad\forall s,t,
    \\
    &G[(z,s),(z,s)]
    =
    P(z|s)
    \quad\forall z,s,
    \\
    &\sum_{m_A,m_B}
      m_Ap_{m_A,m_B}
    \le
    \frac{\bar n}{2\tau},
    \\
    &\sum_{m_A,m_B}
      m_Bp_{m_A,m_B}
    \le
    \frac{\bar n}{2\tau}.
\end{aligned}
    \label{eq:stage2-exact-optimization}
\end{equation}
Here \(e_{\rm ph}(G)\) is the affine A7 functional derived in
Section~\ref{sec:a7-general}.

All constraints are affine in \(G\) and \(p\), except for the PSD constraint
\(G\succeq0\).  Thus the only obstruction is the infinitely many
photon-number variables.

\subsection{Rigorous input truncation}

Let
\[
    S_N=\{(m_A,m_B):m_A+m_B\le N\},
\qquad
    \varepsilon_N
    =
    1-\sum_{S_N}p_{m_A,m_B}.
\]
The separate per-arm energy constraints safely relax to
\begin{equation}
    \sum_{S_N}m_Ap_{m_A,m_B}
    \le\frac{\bar n}{2\tau},
    \qquad
    \sum_{S_N}m_Bp_{m_A,m_B}
    \le\frac{\bar n}{2\tau}.
    \label{eq:stage2-perarm-energy}
\end{equation}
The tail cannot be charged separately to each arm because, e.g.,
\((m_A,m_B)=(0,N+1)\) lies outside \(S_N\).  For the total energy, however,
\(m_A+m_B\ge N+1\) on \(S_N^c\), so
\begin{equation}
    \sum_{S_N}
      (m_A+m_B)p_{m_A,m_B}
    +(N+1)\varepsilon_N
    \le
    \frac{\bar n}{\tau}.
    \label{eq:stage2-combined-energy}
\end{equation}

For each state pair define
\[
    \rho_{st}
    :=
    \max\{|z_A^{st}|,|z_B^{st}|\}.
\]
Then the omitted overlap satisfies
\begin{align}
    \left|
      \sum_{S_N^c}
      p_{m_A,m_B}
      (z_A^{st})^{m_A}(z_B^{st})^{m_B}
    \right|
    &\le
    \sum_{S_N^c}
      p_{m_A,m_B}
      \rho_{st}^{m_A+m_B}
    \nonumber\\
    &\le
    \varepsilon_N\rho_{st}^{N+1}.
    \label{eq:stage2-decaying-tail}
\end{align}
Thus the finite SDP imposes the SOC constraint
\begin{equation}
    \left|
      \sum_zG[(z,s),(z,t)]
      -
      \sum_{S_N}
      p_{m_A,m_B}
      (z_A^{st})^{m_A}(z_B^{st})^{m_B}
    \right|
    \le
    \varepsilon_N\rho_{st}^{N+1}.
    \label{eq:stage2-truncated-soc}
\end{equation}

When both phase differences are nonzero,
\(|z_X^{st}|<1\), so the attenuation produces an exponentially decaying tail
factor.  If one phase difference vanishes then the corresponding
\(|z_X^{st}|=1\), and the uniform bound may revert to the bare
\(\varepsilon_N\).

\subsection{Why Stage 2 becomes computationally impractical}

Pure loss obeys
\[
    \mathbb E[N_X]=\tau\mathbb E[M_X].
\]
Therefore an \(O(1)\) delivered mean requires input photon numbers of order
\[
    M_X=O(\tau^{-1}).
\]
At realistic attenuation such as \(\tau\sim10^{-6}\), a direct Fock grid
would therefore have to reach photon numbers of order \(10^5\)--\(10^6\),
which is computationally prohibitive.

This scaling, rather than a conceptual failure of the SDP, motivates the
continuum reformulation of Stage~3.

For reference, the current finite-Fock numerical values at
\[
    \bar n=10^{-2},\quad L=0,\quad M=4,\quad N_{\rm in}=60
\]
were
\begin{center}
\begin{tabular}{cc}
\toprule
\(\tau=t^2\) & \(e_{\rm ph}\) \\
\midrule
1.0 & 0.2567 \\
0.5 & 0.2293 \\
0.1 & 0.2094 \\
\bottomrule
\end{tabular}
\end{center}
These values are historical numerical benchmarks; they are not the central
rigour statement of this section.

\section{Stage 3: exact rescaling, continuum relaxation, and certified finite representation}
\label{sec:stage3-rigorous}

Stage~3 is best understood as a sequence of representations of the same
mean-energy information.  The starting Stage-2 source is an infinite discrete
distribution \(p_{m_A,m_B}\) over the Fock lattice.  Because the trusted
attenuation is very strong, the monitored post-attenuation energy
\(\bar n\) corresponds to pre-attenuation photon numbers of order
\(1/\tau\).  A direct Fock truncation must therefore become extremely large.

The conceptual chain is
\begin{equation}
\boxed{
p_{m_A,m_B}
\;\xrightarrow[\text{exact}]{u_X=\tau m_X}\;
\nu_\tau
\;\xrightarrow[\text{outer relaxation}]{\text{drop lattice support}}\;
\nu
\;\xrightarrow[\text{outer relaxation}]{\text{cell enclosure}}\;
\{q_j\}_{j=1}^{J}.
}
\label{eq:stage3-chain}
\end{equation}
The same physical mean-energy bound is carried through this entire chain.
Only the last two arrows enlarge the feasible set.

\subsection{Exact push-forward from photon number to delivered energy}

Let \(M_A,M_B\) denote the pre-attenuation photon numbers.  In Stage~2 the
monitored outgoing mean-energy bounds are written in the input variables as
\begin{equation}
    \mathbb E[M_A]\le\frac{\bar n}{2\tau},
    \qquad
    \mathbb E[M_B]\le\frac{\bar n}{2\tau}.
    \label{eq:stage3-input-energy}
\end{equation}
These are first-moment constraints only.  They do \emph{not} impose a hard
photon-number cutoff: arbitrarily large \(m_X\) remain possible with
sufficiently small probability.  Thus the informal phrase
\emph{very large Fock grid} means that a numerical truncation would have to
extend to the \(O(1/\tau)\) scale, and potentially beyond it to capture rare
tails.

Define the rescaled variables
\begin{equation}
    u_A=\tau m_A,
    \qquad
    u_B=\tau m_B.
    \label{eq:stage3-rescaling}
\end{equation}
For a Stage-2 distribution \(p_{m_A,m_B}\), define its exact push-forward
measure
\begin{equation}
    \boxed{
    \nu_\tau
    :=
    \sum_{m_A,m_B\ge0}
      p_{m_A,m_B}
      \delta_{(\tau m_A,\tau m_B)}.
    }
    \label{eq:stage3-lattice-measure}
\end{equation}
It is supported on
\[
    (\tau\mathbb Z_{\ge0})^2.
\]
No approximation has yet been made.

The energy bound transforms exactly:
\begin{align}
    \int u_A\,d\nu_\tau
    &=
    \tau\sum_{m_A,m_B}m_Ap_{m_A,m_B}
    \le\frac{\bar n}{2},
    \\
    \int u_B\,d\nu_\tau
    &=
    \tau\sum_{m_A,m_B}m_Bp_{m_A,m_B}
    \le\frac{\bar n}{2}.
    \label{eq:stage3-perarm-energy}
\end{align}
Hence
\begin{equation}
    \int(u_A+u_B)\,d\nu_\tau
    \le\bar n.
    \label{eq:stage3-energy}
\end{equation}
The factor \(1/\tau\) has disappeared because the natural Stage-3 variables
are the delivered-energy variables \(u_X=\tau m_X\).

\subsection{Exact finite-\texorpdfstring{\(\tau\)}{tau} kernel on the rescaled variables}

For a state pair \(s,t\), define
\[
    \kappa_X^{st}
    :=
    1-\tau+\tau e^{i\Delta\phi_X^{st}}
    =
    1-\tau a_X^{st},
    \qquad
    a_X^{st}:=1-e^{i\Delta\phi_X^{st}}.
\]
On the physical lattice,
\[
    (\kappa_A^{st})^{m_A}(\kappa_B^{st})^{m_B}
\]
is exactly the Stage-2 binomial-thinning kernel.

For the strong-attenuation regime \(0<\tau<1/2\),
\(\kappa_X^{st}\) lies in the open right half-plane, so the principal logarithm
is unambiguous.  Define
\begin{equation}
    b_X^{st}
    :=
    \frac{1}{\tau}\operatorname{Log}\kappa_X^{st},
    \qquad
    \boxed{
    K_\tau^{st}(u_A,u_B)
    :=
    \exp\!\left[
      b_A^{st}u_A+b_B^{st}u_B
    \right].
    }
    \label{eq:stage3-finite-kernel}
\end{equation}
When \(u_X=\tau m_X\), this reduces exactly to
\[
    K_\tau^{st}(\tau m_A,\tau m_B)
    =
    (\kappa_A^{st})^{m_A}(\kappa_B^{st})^{m_B}.
\]
Therefore the exact physical overlap is
\begin{equation}
    \boxed{
    \Lambda^{(\tau)}_{st}
    =
    \int
      K_\tau^{st}(u_A,u_B)
      \,d\nu_\tau(u_A,u_B).
    }
    \label{eq:stage3-exact-integral}
\end{equation}

Because \(|\kappa_X^{st}|\le1\), one has
\[
    \Re b_X^{st}\le0,
    \qquad
    |K_\tau^{st}(u)|\le1
    \quad
    (u_A,u_B\ge0).
\]

\subsection{Dropping the lattice while keeping the energy bound}

Stage~3 now drops only the support condition
\[
    \operatorname{supp}\nu_\tau
    \subset
    (\tau\mathbb Z_{\ge0})^2
\]
and allows an arbitrary positive probability measure
\[
    \nu
    \quad\text{on}\quad
    \mathbb R_{\ge0}^2,
\]
while retaining the same first-moment constraints
\begin{equation}
    \boxed{
    \int u_A\,d\nu\le\frac{\bar n}{2},
    \qquad
    \int u_B\,d\nu\le\frac{\bar n}{2}.
    }
    \label{eq:stage3-continuum-energy}
\end{equation}

Thus
\begin{equation}
    \boxed{
    \{\text{physical lattice measures }\nu_\tau\}
    \subseteq
    \{\text{continuum measures }\nu
       \text{ satisfying \eqref{eq:stage3-continuum-energy}}\}.
    }
    \label{eq:stage3-lattice-inclusion}
\end{equation}
This is a genuine outer relaxation: the energy budget is unchanged, but Eve
is no longer required to place her probability mass on the
\(\tau\)-spaced lattice.

The exact finite-\(\tau\) kernel can be retained after this relaxation.  The
formal continuum problem is therefore
\begin{equation}
\begin{aligned}
    \text{maximize}_{G,\nu}\quad&
    e_{\rm ph}(G)
    \\
    \text{subject to}\quad&
    G\succeq0,
    \\
    &\nu\ge0,
    \qquad
    \int d\nu=1,
    \\
    &\int u_A\,d\nu\le\frac{\bar n}{2},
    \qquad
    \int u_B\,d\nu\le\frac{\bar n}{2},
    \\
    &\sum_zG[(z,s),(z,t)]
      =
      \int K_\tau^{st}(u)\,d\nu(u)
      \quad\forall s,t,
    \\
    &G[(z,s),(z,s)]
    =
    P(z|s)
    \quad\forall z,s.
\end{aligned}
    \label{eq:stage3-measure-program}
\end{equation}
The equality in Eq.~\eqref{eq:stage3-measure-program} is simply overlap
conservation:
\[
    \underbrace{\int K_\tau^{st}\,d\nu}_{\text{input overlap}}
    =
    \underbrace{\Lambda_{st}}_{\text{same overlap}}
    =
    \underbrace{\sum_zG[(z,s),(z,t)]}_{\text{after Eve's instrument}}.
\]
The measure-valued program is convex, but infinite-dimensional.

\subsection{The limiting kernel \(K_0\): useful interpretation, not required numerically}

At fixed \(u_X\),
\[
    b_X^{st}
    =
    -a_X^{st}+O(\tau),
\]
so
\begin{equation}
    K_\tau^{st}(u_A,u_B)
    \longrightarrow
    K_0^{st}(u_A,u_B)
    :=
    \exp[
      -u_Aa_A^{st}
      -u_Ba_B^{st}
    ].
    \label{eq:stage3-limit-kernel}
\end{equation}
The limiting kernel has the mixed-Poisson interpretation
\[
    N_X|U_X=u_X
    \sim
    \operatorname{Poisson}(u_X),
\]
because
\[
    \mathbb E[e^{i\theta N_X}|U_X=u_X]
    =
    e^{-u_X(1-e^{i\theta})}.
\]
Thus the continuum model is naturally interpreted as
\[
    (U_A,U_B)\sim\nu,
\]
followed conditionally by independent Poisson photon numbers in the two arms.
The variables \(U_A,U_B\) may remain arbitrarily correlated.

A coherent-state photon-number law corresponds to a delta measure in \(u\).
Strong attenuation does \emph{not} force an arbitrary source to become
coherent; an arbitrary mixing measure \(\nu\) remains.

For the certified numerical implementation there is no need to replace
\(K_\tau\) by \(K_0\): Eq.~\eqref{eq:stage3-finite-kernel} can be used
directly.  If \(K_0\) is used for analytic estimates, the bound
\begin{equation}
    |K_\tau^{st}-K_0^{st}|
    \le
    4\tau(u_A+u_B)
    e^{4\tau(u_A+u_B)}
    \label{eq:stage3-pointwise-error}
\end{equation}
provides an explicit finite-\(\tau\) error.  On a region
\(u_A+u_B\le U\),
\begin{equation}
    \left|
      \int K_\tau^{st}\,d\nu
      -
      \int K_0^{st}\,d\nu
    \right|
    \le
    4\tau\bar n\,e^{4\tau U}.
    \label{eq:stage3-integrated-error}
\end{equation}
This \(K_0\)-replacement is therefore optional, not part of the preferred
certified solver.

\subsection{Only finitely many moments of \(\nu\) are visible to the SDP}

Although \(\nu\) is infinite-dimensional, the QKD SDP sees it only through
normalization, first moments, and the real and imaginary parts of finitely many
kernel integrals
\[
    \int K_\tau^{st}\,d\nu.
\]
On a compact domain \(D\), define
\[
    F(u)
    =
    \left(
      1,u_A,u_B,
      \{\Re K_\tau^{st}(u),\Im K_\tau^{st}(u)\}_{s<t}
    \right).
\]
Then
\[
    y=\int_DF(u)\,d\nu(u)
\]
lies in the convex hull of \(\{F(u):u\in D\}\).  Standard convex-hull results
therefore guarantee the existence of finite atomic representations of the
visible moment vector.  However, the required atom locations are unknown
continuous variables.  Optimizing their locations jointly with their weights
does not preserve the simple finite SDP structure.  This is why a certified
cell enclosure is useful.

\subsection{Continuum tail truncation}

For computation choose
\[
    D_U
    :=
    \{(u_A,u_B):u_A+u_B\le U\}
\]
and define
\[
    \varepsilon_U
    :=
    \nu(D_U^c).
\]
Since every point outside \(D_U\) has \(u_A+u_B>U\),
\begin{equation}
    \int_{D_U}(u_A+u_B)d\nu
    +
    U\varepsilon_U
    \le
    \bar n.
    \label{eq:stage3-tail-energy}
\end{equation}
Because \(|K_\tau^{st}|\le1\),
\begin{equation}
    \left|
      \int_{D_U^c}K_\tau^{st}\,d\nu
    \right|
    \le
    \varepsilon_U.
    \label{eq:stage3-tail-overlap}
\end{equation}
The separate in-domain bounds
\[
    \int_{D_U}u_A\,d\nu\le\frac{\bar n}{2},
    \qquad
    \int_{D_U}u_B\,d\nu\le\frac{\bar n}{2}
\]
also remain valid by dropping the nonnegative tail contribution.  As before,
the same \(U\varepsilon_U\) cannot be charged separately to both arms.

\subsection{Why a fixed support grid is an inner approximation}

A fixed-grid discretization assumes
\[
    \nu
    =
    \sum_{j=1}^Jq_j\delta_{c_j}
\]
for predetermined support points \(c_j\).  It therefore restricts Eve:
\[
    \{\text{grid-supported measures}\}
    \subset
    \{\text{all continuum measures}\}.
\]
Hence
\begin{equation}
    e_{\rm ph}^{\rm grid}
    \le
    e_{\rm ph}^{\rm continuum},
    \label{eq:stage3-grid-inner}
\end{equation}
which is the wrong direction for a security upper bound.  A very dense grid
may be a useful numerical diagnostic, but density alone does not certify it.

Figure~\ref{fig:stage3-cell-relaxation} summarizes the corrected chain.  The
first arrow carries the energy bound exactly from the large Fock lattice to
the rescaled variables.  The second arrow drops only the lattice support.
The final finite variables are cell masses, not point masses.

\begin{figure}[t]
\centering
\begin{tikzpicture}[
    >=Latex,
    font=\small,
    axis/.style={->,thick},
    domain/.style={draw=blue!65!black,thick,fill=blue!5},
    cell/.style={draw=black!45,thin},
    hilite/.style={draw=green!50!black,thick,fill=green!18},
    note/.style={rounded corners,draw=black!40,fill=white,inner sep=4pt,align=center}
]

% ------------------------------------------------------------------
% Left panel: continuum measure
% ------------------------------------------------------------------
\node[font=\bfseries] at (2.0,4.65) {Continuum measure};

\draw[axis] (0,0) -- (4.45,0) node[right] {$u_A$};
\draw[axis] (0,0) -- (0,4.45) node[above] {$u_B$};
\draw[domain] (0,0) -- (0,4) -- (4,0) -- cycle;
\node[anchor=east] at (-0.08,4) {$U$};
\node[anchor=north] at (4,-0.08) {$U$};
\node[anchor=north east] at (-0.05,-0.05) {$0$};

% A schematic non-atomic measure cloud
\fill[blue!35,opacity=.28] (1.15,1.15) circle (0.72);
\fill[cyan!55,opacity=.24] (1.55,0.82) circle (0.58);
\fill[orange!65,opacity=.20] (0.85,1.55) circle (0.48);
\fill[blue!55,opacity=.18] (2.35,0.72) circle (0.42);
\node[align=center] at (1.25,2.30)
    {$\nu$ arbitrary on\\[-1mm] $D_U=\{u_A+u_B\le U\}$};

\node[note,text width=4.15cm] at (2.05,-1.05)
    {$\displaystyle
      \Lambda_{st}
      =\int_{D_U} K_{\tau}^{st}(u)\,d\nu(u)
      +r^{\rm tail}_{st}$\\
     $\displaystyle |r^{\rm tail}_{st}|\le\varepsilon$};

% Transition arrow
\draw[->,very thick,gray!65] (4.85,2.0) -- (6.10,2.0);
\node[align=center,text=black!70] at (5.48,2.55)
    {\footnotesize partition\\[-1mm]\footnotesize into cells};

% ------------------------------------------------------------------
% Right panel: cell masses
% ------------------------------------------------------------------
\begin{scope}[xshift=6.55cm]
\node[font=\bfseries] at (2.0,4.65) {Finite cell variables};

\draw[axis] (0,0) -- (4.45,0) node[right] {$u_A$};
\draw[axis] (0,0) -- (0,4.45) node[above] {$u_B$};
\draw[domain] (0,0) -- (0,4) -- (4,0) -- cycle;
\node[anchor=east] at (-0.08,4) {$U$};
\node[anchor=north] at (4,-0.08) {$U$};
\node[anchor=north east] at (-0.05,-0.05) {$0$};

% Grid, clipped to triangular domain
\begin{scope}
\clip (0,0) -- (0,4) -- (4,0) -- cycle;
\foreach \x in {0.8,1.6,2.4,3.2}
    \draw[cell] (\x,0) -- (\x,4);
\foreach \y in {0.8,1.6,2.4,3.2}
    \draw[cell] (0,\y) -- (4,\y);
\end{scope}

% Highlight one complete cell
\filldraw[hilite] (0.8,0.8) rectangle (1.6,1.6);
\node at (1.04,1.40) {$C_j$};
\fill[black] (1.20,1.20) circle (1.5pt);
\node[anchor=west] at (1.27,1.19) {$c_j$};

% Show that the actual mass need not sit at c_j
\fill[green!45!black] (0.94,0.98) circle (1.1pt);
\fill[green!45!black] (1.48,1.08) circle (1.1pt);
\fill[green!45!black] (1.36,1.48) circle (1.1pt);

\node[note,text width=3.15cm] (massnote) at (2.85,2.85)
    {$q_j:=\nu(C_j)$\\[1mm]
     mass may lie anywhere in $C_j$};
\draw[->,green!45!black,thick]
    (massnote.south west) .. controls (2.25,2.25) and (1.85,1.75) .. (1.53,1.47);

\node[note,text width=4.35cm] at (2.05,-1.05)
    {$\displaystyle
      \int_{C_j}K_{\tau}^{st}(u)\,d\nu(u)
      =q_jK_{\tau}^{st}(c_j)+r_{st,j}$\\
     $\displaystyle |r_{st,j}|\le q_j\delta_{st,j}$};
\end{scope}

% ------------------------------------------------------------------
% Certified global constraint
% ------------------------------------------------------------------
\node[
    rounded corners,
    draw=orange!65!black,
    fill=orange!6,
    inner sep=7pt,
    text width=11.2cm,
    align=center
] at (5.30,-2.55)
{
\[
\left|
\sum_z G[(z,s),(z,t)]
-
\sum_j q_j K_{\tau}^{st}(c_j)
\right|
\le
\sum_j q_j\delta_{st,j}
+\varepsilon ,
\qquad
\delta_{st,j}\ge
\sup_{u\in C_j}
\left|K_{\tau}^{st}(u)-K_{\tau}^{st}(c_j)\right|.
\]
};

\node[
    rounded corners,
    draw=green!45!black,
    fill=green!5,
    inner sep=5pt,
    text width=10.7cm,
    align=center
] at (5.30,-3.95)
{
Every continuum measure maps to feasible cell variables:
\[
q_j=\nu(C_j),
\qquad
\varepsilon=\nu(D_U^c),
\qquad
\mathcal F_{\rm continuum}\subseteq\mathcal F_{\rm cell}.
\]
The representative points \(c_j\) are only used to evaluate the kernel;
they are \emph{not} imposed support points of \(\nu\).
};

\end{tikzpicture}
\caption{Schematic of the certified continuum-to-cell relaxation in Stage~3.
The continuum measure is replaced by cell masses \(q_j=\nu(C_j)\), not by
point masses at the representatives \(c_j\).  The variables
\(\delta_{st,j}\) bound the possible kernel variation within each cell, while
\(\varepsilon\) bounds the omitted tail.  Consequently every feasible
continuum measure produces a feasible point of the finite conic program,
which is the required outer-approximation direction.  The figure is written
with the exact finite-\(\tau\) kernel \(K_\tau\); the same construction applies
to the limiting kernel \(K_0\).}
\label{fig:stage3-cell-relaxation}
\end{figure}


\subsection{Certified lifted cell relaxation}

Partition the represented domain as
\[
    D_U=\bigcup_{j=1}^{J}C_j.
\]
For each cell retain not only its probability mass
\begin{equation}
    q_j:=\nu(C_j),
    \label{eq:stage3-cell-mass}
\end{equation}
but also its two first moments
\begin{equation}
    x_{A,j}
    :=
    \int_{C_j}u_A\,d\nu,
    \qquad
    x_{B,j}
    :=
    \int_{C_j}u_B\,d\nu.
    \label{eq:stage3-cell-first-moments}
\end{equation}
Thus \(q_j\) says how much mass lies somewhere in \(C_j\), while
\(x_{A,j},x_{B,j}\) retain the mean location of that mass inside the cell.
No point-support assumption is made.

Let
\[
    \underline u_{X,j}:=\inf_{u\in C_j}u_X,
    \qquad
    \overline u_{X,j}:=\sup_{u\in C_j}u_X,
\]
and similarly
\[
    \underline s_j:=\inf_{u\in C_j}(u_A+u_B),
    \qquad
    \overline s_j:=\sup_{u\in C_j}(u_A+u_B).
\]
Every genuine continuum measure satisfies the perspective constraints
\begin{align}
    \underline u_{A,j}q_j
    &\le x_{A,j}\le
    \overline u_{A,j}q_j,
    \\
    \underline u_{B,j}q_j
    &\le x_{B,j}\le
    \overline u_{B,j}q_j,
    \\
    \underline s_jq_j
    &\le x_{A,j}+x_{B,j}\le
    \overline s_jq_j.
    \label{eq:stage3-cell-perspective}
\end{align}
For the clipped rectangular cells used numerically these inequalities encode
the fact that the conditional barycentre of the mass in \(C_j\) lies inside
the same convex cell.

\subsubsection{First-order kernel lifting and second-order remainder}

Choose a representative point
\[
    c_j=(c_{A,j},c_{B,j})\in C_j
\]
only as a linearisation point.  Recall that
\[
    K_\tau^{st}(u)
    =
    \exp[
      b_A^{st}u_A+b_B^{st}u_B
    ].
\]
For \(u\in C_j\), write
\[
    \Delta u_X=u_X-c_{X,j},
    \qquad
    z_{st,j}(u)
    :=
    b_A^{st}\Delta u_A
    +
    b_B^{st}\Delta u_B.
\]
Then
\[
    K_\tau^{st}(u)
    =
    K_\tau^{st}(c_j)e^{z_{st,j}(u)}.
\]
Using
\[
    e^z=1+z+R_2(z),
    \qquad
    |R_2(z)|
    \le
    \frac{|z|^2}{2}e^{|z|},
\]
the integral over one cell can be written as
\begin{align}
    \int_{C_j}K_\tau^{st}(u)\,d\nu(u)
    &=
    K_\tau^{st}(c_j)
    \Big[
      q_j
      +
      b_A^{st}(x_{A,j}-c_{A,j}q_j)
      \nonumber\\
    &\hspace{29mm}
      +
      b_B^{st}(x_{B,j}-c_{B,j}q_j)
    \Big]
    +
    r_{st,j}.
    \label{eq:stage3-lifted-cell-decomposition}
\end{align}
The first-order term is therefore represented \emph{exactly} by the lifted
variables \(q_j,x_{A,j},x_{B,j}\).

Define
\begin{equation}
    \rho_{st,j}
    :=
    \sup_{u\in C_j}
    \left|
      b_A^{st}(u_A-c_{A,j})
      +
      b_B^{st}(u_B-c_{B,j})
    \right|.
    \label{eq:stage3-cell-bound}
\end{equation}
Since the absolute value of an affine function is convex, for a polygonal
cell the supremum is attained at a vertex and is easy to compute exactly.
The Taylor remainder obeys
\begin{equation}
    |r_{st,j}|
    \le
    q_j\,\eta_{st,j},
    \qquad
    \eta_{st,j}
    :=
    \frac{\rho_{st,j}^2}{2}e^{\rho_{st,j}}.
    \label{eq:stage3-second-order-remainder}
\end{equation}
A second independent bound follows from
\(|K_\tau(u)|,|K_\tau(c_j)|\le1\):
\[
    |r_{st,j}|
    \le
    q_j(2+\rho_{st,j}).
\]
The implementation uses the smaller of these two certified bounds.

Define the lifted finite approximation
\begin{align}
    \widetilde\Lambda_{st}^{\rm cell}
    :=
    \sum_j
    K_\tau^{st}(c_j)
    \Big[
      q_j
      &+
      b_A^{st}(x_{A,j}-c_{A,j}q_j)
      \nonumber\\
      &+
      b_B^{st}(x_{B,j}-c_{B,j}q_j)
    \Big].
    \label{eq:stage3-lifted-overlap}
\end{align}
Including the unknown tail outside \(D_U\), whose overlap contribution has
modulus at most \(\varepsilon_U\), gives the certified SOC constraint
\begin{equation}
\boxed{
    \left|
      \sum_zG[(z,s),(z,t)]
      -
      \widetilde\Lambda_{st}^{\rm cell}
    \right|
    \le
    \sum_jq_j\eta_{st,j}
    +
    \varepsilon_U.
}
    \label{eq:stage3-certified-soc}
\end{equation}
The important improvement over the zeroth-order cell enclosure is that the
within-cell error is now second order in the cell diameter rather than first
order.

\subsection{How the energy bound reaches the lifted cell variables}

The lifted moments also remove most of the looseness in the energy
constraint.  By definition,
\[
    \sum_jx_{A,j}
    =
    \int_{D_U}u_A\,d\nu,
    \qquad
    \sum_jx_{B,j}
    =
    \int_{D_U}u_B\,d\nu.
\]
Hence every feasible continuum measure satisfies
\begin{equation}
    \boxed{
    \sum_jx_{A,j}\le\frac{\bar n}{2},
    \qquad
    \sum_jx_{B,j}\le\frac{\bar n}{2}.
    }
    \label{eq:stage3-cell-perarm-energy}
\end{equation}
Together with
\[
    \sum_jq_j+\varepsilon_U=1,
\]
the triangular tail gives
\begin{equation}
    \boxed{
    \sum_j(x_{A,j}+x_{B,j})
    +
    U\varepsilon_U
    \le
    \bar n.
    }
    \label{eq:stage3-cell-total-energy}
\end{equation}

Thus the same physical energy statement is carried through the entire chain:
\begin{equation}
\boxed{
\mathbb E[M_X]\le\frac{\bar n}{2\tau}
\;\Longrightarrow\;
\int u_X\,d\nu_\tau\le\frac{\bar n}{2}
\;\Longrightarrow\;
\int u_X\,d\nu\le\frac{\bar n}{2}
\;\Longrightarrow\;
\sum_jx_{X,j}\le\frac{\bar n}{2}.
}
\label{eq:stage3-energy-chain}
\end{equation}
The first arrow is an exact change of variables.  The middle step enlarges
only the support of the measure.  The final step does not approximate the
in-domain first moment: it introduces variables that are precisely those cell
moments.  Only the unrepresented tail is relaxed through the lower charge
\(U\varepsilon_U\).

\subsection{Proof of the lifted outer-approximation direction}

Take any feasible continuum measure \(\nu\).  Define
\[
    q_j=\nu(C_j),
    \qquad
    x_{A,j}=\int_{C_j}u_A\,d\nu,
    \qquad
    x_{B,j}=\int_{C_j}u_B\,d\nu,
    \qquad
    \varepsilon_U=\nu(D_U^c).
\]
The perspective constraints
Eq.~\eqref{eq:stage3-cell-perspective} follow immediately because each
conditional cell barycentre lies in \(C_j\).  The energy constraints
Eqs.~\eqref{eq:stage3-cell-perarm-energy}--\eqref{eq:stage3-cell-total-energy}
follow from the original first-moment bounds.

For every state pair,
Eq.~\eqref{eq:stage3-lifted-cell-decomposition} and the certified remainder
bound imply
\[
    \int K_\tau^{st}\,d\nu
    =
    \widetilde\Lambda_{st}^{\rm cell}
    +
    \sum_jr_{st,j}
    +
    r_{st}^{\rm tail},
\]
where
\[
    |r_{st,j}|
    \le
    q_j\eta_{st,j},
    \qquad
    |r_{st}^{\rm tail}|
    \le
    \varepsilon_U.
\]
Thus every continuum-feasible point maps to a feasible lifted finite point:
\begin{equation}
    \boxed{
    \mathcal F_{\rm physical}
    \subseteq
    \mathcal F_{\rm continuum}
    \subseteq
    \mathcal F_{\rm lifted\ cell}.
    }
    \label{eq:stage3-feasible-chain}
\end{equation}
Consequently,
\begin{equation}
    \boxed{
    e_{\rm ph}^{\rm physical}
    \le
    e_{\rm ph}^{\rm continuum}
    \le
    e_{\rm ph}^{\rm lifted\ cell}.
    }
    \label{eq:stage3-objective-chain}
\end{equation}

\begin{remark}[Cell refinement and monotonicity]
The second-order remainders shrink rapidly with cell diameter.  Nevertheless,
independently constructed relaxations at two resolutions are not
automatically nested.  A strictly monotone hierarchy can be enforced by
retaining the parent-cell constraints when child cells are introduced.
\end{remark}

\subsection{Alternative computational route: adaptive exchange}

A potentially more efficient alternative is a semi-infinite exchange method:

\begin{enumerate}
    \item solve a finite SDP on a current set of support points;
    \item use the dual solution to construct a continuous violation function
    over \(u\in D_U\);
    \item globally minimize that function;
    \item if a violating point exists, add it and repeat.
\end{enumerate}

Because \(u=(u_A,u_B)\) is only two-dimensional, this may eventually be
attractive.  Rigour would require a certified global solution of the
continuous separation problem.  The cell-enclosure route remains the
preferred first implementation.

\subsection{Mixed-Poisson moments and concentration freedom}

For one arm in the limiting model,
\[
    N|U=u\sim\operatorname{Poisson}(u).
\]
Hence
\[
    \mathbb E[N]
    =
    \mathbb E[U],
    \qquad
    \mathbb E[N(N-1)]
    =
    \mathbb E[U^2],
\]
and
\begin{equation}
    g^{(2)}
    =
    1+
    \frac{\operatorname{Var}(U)}
         {\mathbb E[U]^2}.
    \label{eq:stage3-mixed-poisson-g2}
\end{equation}
A first-moment constraint fixes \(\mathbb E[U]\) but leaves
\(\operatorname{Var}(U)\) unbounded.  This is the continuum expression of the
rare-bright-pulse freedom already visible in Stage~1.

\subsection{Status of the present Stage-3 numerics}

At \(\tau=10^{-6}\), \(M=4\), and \(L=0\), the earlier fixed-support scan was
largest near
\[
    \bar n\simeq0.1,
    \qquad
    e_{\rm ph}\simeq0.235\text{--}0.238,
    \qquad
    R\simeq7\times10^{-3}.
\]
Those numbers remain useful historical diagnostics, but the fixed-support
calculation is an inner approximation of the continuum problem and is not a
security certificate.

The code \texttt{phase/pnp-stage-3.py} now implements the certified lifted
cell relaxation of Eq.~\eqref{eq:stage3-certified-soc} using the exact
finite-\(\tau\) kernel \(K_\tau\), cell masses
\(q_j=\nu(C_j)\), cell first moments \(x_{A,j},x_{B,j}\), a certified
second-order kernel remainder, and a continuum-tail variable.
The remaining numerical task is to benchmark this certified outer bound
against the older fixed-support values, study cell refinement, and control
solver tolerances.  Until those scans are rerun, the numerical values quoted
above should still be read only as diagnostics.

\section{Concentration attack: rigorous inequality and heuristic saturation}
\label{sec:concentration-rigorous}

Consider the two-point family
\begin{equation}
    \nu
    =
    (1-p)\delta_0+p\delta_{u_0}.
    \label{eq:stage3-concentration-family}
\end{equation}
Let \(q(u)\in[0,1]\) be the conditional probability of a pass announcement.

If vacuum gives no pass, \(q(0)=0\), then
\[
    P_{\rm pass}
    =
    p\,q(u_0)
    \le
    p.
\]
Thus \(p\ge P_{\rm pass}\). With the per-arm mean budget
\(\mu=\bar n/2\),
\[
    pu_0\le\mu,
\]
so
\begin{equation}
    u_0
    \le
    \frac{\bar n/2}{P_{\rm pass}}.
    \label{eq:stage3-concentration-exact}
\end{equation}
With vacuum pass probability \(q_0>0\),
\[
    P_{\rm pass}
    \le
    q_0+p(1-q_0),
\]
which gives
\begin{equation}
    p
    \ge
    \frac{P_{\rm pass}-q_0}{1-q_0},
    \qquad
    u_0
    \le
    \frac{\bar n}{2}
    \frac{1-q_0}{P_{\rm pass}-q_0},
    \label{eq:stage3-concentration-dark}
\end{equation}
provided \(P_{\rm pass}>q_0\).

These are rigorous upper bounds for the two-point family. The claim that Eve
can nearly attain them requires \(q(u_0)=O(1)\) for bright pulses. In that
regime \(p\approx P_{\rm pass}\).

For the honest weak-signal model,
\[
    P_{\rm pass}
    \simeq
    c\,\eta\,\bar n
\]
with a convention-dependent \(c=O(1)\). Then
\begin{equation}
    u_{0,\max}
    =
    O(\eta^{-1}).
    \label{eq:stage3-one-over-eta}
\end{equation}
The cancellation of \(\bar n\) is therefore a scaling statement derived from
the honest pass model.

For distinct phases,
\begin{equation}
    |K_0(u_0;\Delta\phi)|
    =
    e^{-u_0(1-\cos\Delta\phi)},
    \label{eq:stage3-bright-overlap}
\end{equation}
so bright components become exponentially distinguishable. This explains why
the optimized phase-error SDP moves toward \(1/2\) with distance. The logical
step from vanishing pairwise overlap to \(e_{\rm ph}\to1/2\), however, is not
a standalone theorem; it is an interpretation of the SDP optimum.

\section{Stage 4: abstract conditional upper-tail monitoring}
\label{sec:stage4-rigorous}

\subsection{Final scope of Stage 4}

Stage~4 is the final \emph{abstract monitor} stage.  It deliberately does not
model a physical APN or PNA detector, beam splitter, detector noise, or
measurement back-action.  Instead it asks a narrower question:

\begin{quote}
If a trusted local monitor can certify that the monitor-accepted ensemble has
only a small high-energy tail, is that information sufficient to recover a
useful phase-error bound against an otherwise arbitrary untrusted source?
\end{quote}

The answer of the present numerical and analytic study is yes.

This should be contrasted with Stage~3.  Stage~3 carries only first-moment
information,
\begin{equation}
    \mathbb E[U_X]\le\frac{\bar n}{2},
    \label{eq:stage4-stage3-mean}
\end{equation}
which is APN-like information and permits arbitrarily remote tails.  Stage~4
replaces that information by a conditional upper-tail statement.  No lower
energy threshold is required: vacuum and weak accepted pulses are harmless
from the concentration-attack point of view, whereas rare arbitrarily bright
accepted pulses are the dangerous degree of freedom.

The variable used in the current Stage-4 implementation is still the Stage-3
effective delivered-energy variable
\begin{equation}
    U_X=\tau M_X,
    \label{eq:stage4-u-definition}
\end{equation}
where \(M_X\) is the photon number before the trusted pure-loss map of
transmissivity \(\tau\).  Thus \(U_X\) is a rescaled source variable (and the
conditional mean output photon number for a fixed input Fock number), not the
actual realised output photon number of a particular round.

\subsection{Abstract monitor interface}

Let \(A_{\rm mon}\) denote acceptance by the trusted local monitor.  Define
\begin{equation}
    W_+
    :=
    [0,u_{A,+}]
    \times
    [0,u_{B,+}].
    \label{eq:stage4-good-region}
\end{equation}
The Stage-4 security interface is
\begin{equation}
    \boxed{
    \Pr\!\left[
      U_A>u_{A,+}
      \ \text{or}\
      U_B>u_{B,+}
      \,\middle|\,
      A_{\rm mon}
    \right]
    \le
    \epsilon_{\rm mon}.
    }
    \label{eq:stage4-monitor-interface}
\end{equation}
The parameter \(\epsilon_{\rm mon}\) is therefore a leakage probability for
the \emph{accepted} ensemble.  In the current code
\begin{equation}
    \epsilon_{\rm mon}=10^{-6}
    \label{eq:stage4-current-eps}
\end{equation}
is only a placeholder used for numerical exploration; it is not yet a
detector-derived experimental claim.

Equation~\eqref{eq:stage4-monitor-interface} is stronger in the relevant
direction than a mean-energy bound.  In the good accepted component,
\(U_X\le u_{X,+}\) independently of channel distance and independently of the
central pass probability.  Hence the Stage-3 concentration scale
\[
    u_0\sim\frac{\bar n/2}{P_{\rm pass}}
\]
cannot grow without bound.

\subsection{Conditioning and the common phase orbit}

The monitor is assumed to act before the secret phase and bit choices.  Let
\(\mathcal M_{\rm acc}\) denote its accepted branch.  For an arbitrary
Eve-supplied state \(\rho\),
\begin{equation}
    \sigma_{\rm acc}
    =
    \frac{\mathcal M_{\rm acc}(\rho)}
         {P_{\rm mon}},
    \qquad
    P_{\rm mon}
    =
    \operatorname{Tr}\mathcal M_{\rm acc}(\rho).
    \label{eq:stage4-conditioned-state}
\end{equation}
The accepted state may be mixed, correlated between Alice and Bob, and
entangled with Eve.  None of this obstructs the phase-orbit construction.
After all basis-independent trusted operations have been completed, the
secret encoding itself generates
\begin{equation}
    \rho_s
    =
    U_s\sigma_0U_s^\dagger,
    \label{eq:stage4-common-orbit}
\end{equation}
where \(\sigma_0\) is the accepted reference state at the point of encoding.
Thus the common unitary orbit does \emph{not} have to be restored by
attenuation.  It is automatic as long as monitoring/attenuation are performed
before the phase label \(s\) is applied.

Choosing one purification \(|\chi_0\rangle\) of \(\sigma_0\) and defining
\[
    |\chi_s\rangle
    =
    (U_s\otimes I_R)|\chi_0\rangle
\]
gives
\begin{equation}
    \Lambda_{st}
    =
    \langle\chi_s|\chi_t\rangle
    =
    \operatorname{Tr}
    [\sigma_0U_s^\dagger U_t],
    \label{eq:stage4-purified-overlap}
\end{equation}
which is the overlap required by the Gram SDP.

In the present Stage-4 abstraction the relation between the unknown accepted
source and these overlaps is represented through the same exact finite-\(\tau\)
kernel as Stage~3:
\begin{equation}
    K_\tau^{st}(u_A,u_B)
    =
    \exp[
      b_A^{st}u_A+b_B^{st}u_B
    ],
    \qquad
    b_X^{st}
    =
    \frac{1}{\tau}
    \operatorname{Log}
    \left(
      1-\tau+\tau e^{i\Delta\phi_X^{st}}
    \right).
    \label{eq:stage4-kernel}
\end{equation}

\subsection{Good--bad decomposition}

Let \(\nu_{\rm acc}\) be the accepted effective-energy measure.  Write
\begin{equation}
    \boxed{
    \nu_{\rm acc}
    =
    \nu_{\rm good}
    +
    \nu_{\rm bad},
    }
    \label{eq:stage4-good-bad}
\end{equation}
with
\[
    \operatorname{supp}\nu_{\rm good}
    \subseteq W_+,
\]
and define
\[
    \beta
    :=
    \int d\nu_{\rm bad}.
\]
The monitor statement implies
\begin{equation}
    \boxed{
    0\le\beta\le\epsilon_{\rm mon},
    \qquad
    \int d\nu_{\rm good}+\beta=1.
    }
    \label{eq:stage4-beta-bound}
\end{equation}

For every signal pair,
\begin{equation}
    \Lambda_{st}^{\rm acc}
    =
    \int_{W_+}
      K_\tau^{st}(u)\,d\nu_{\rm good}(u)
    +
    r_{st}^{\rm mon},
    \label{eq:stage4-overlap-decomposition}
\end{equation}
where the bad component is left completely arbitrary.  Since
\(|K_\tau^{st}(u)|\le1\),
\begin{equation}
    \boxed{
    |r_{st}^{\rm mon}|
    \le
    \beta
    \le
    \epsilon_{\rm mon}.
    }
    \label{eq:stage4-monitor-overlap-slack}
\end{equation}

This is the essential Stage-4 reduction: the unbounded continuum tail of
Stage~3 is replaced by a compact good region plus one explicit bad-mass
parameter.

\subsection{No generic conditional mean constraint}

An unconditional APN statement such as
\[
    \mathbb E[U_X]\le\frac{\bar n}{2}
\]
does not imply
\[
    \mathbb E[U_X\mid A_{\rm mon}]
    \le\frac{\bar n}{2}.
\]
Conditioning can change the mean.  Therefore the generic Stage-4 SDP does
\emph{not} impose the Stage-3 mean constraint on \(\nu_{\rm acc}\).  If an
experiment separately certifies a conditional mean bound, it can be added,
but it is not part of the present security claim.

For numerical scans it is convenient to parameterise
\begin{equation}
    u_{A,+}=u_{B,+}
    =
    c\,\frac{\bar n_{\rm hon}}{2},
    \label{eq:stage4-cap-parameterization}
\end{equation}
where \(\bar n_{\rm hon}\) is a nominal honest outgoing brightness used only
as a reference scale.  Equation~\eqref{eq:stage4-cap-parameterization} is
\emph{not} a security identity and \(\bar n_{\rm hon}\) is not an optimized
moment of Eve's Stage-4 distribution.

\subsection{Certified finite-dimensional relaxation}

Partition the compact good region into cells
\[
    W_+
    =
    \bigcup_{j=1}^{J}C_j.
\]
For each cell retain the good mass and its first moments,
\begin{equation}
    q_j
    :=
    \nu_{\rm good}(C_j),
    \qquad
    x_{A,j}
    :=
    \int_{C_j}u_A\,d\nu_{\rm good},
    \qquad
    x_{B,j}
    :=
    \int_{C_j}u_B\,d\nu_{\rm good}.
    \label{eq:stage4-cell-vars}
\end{equation}
They obey the usual perspective constraints
\begin{align}
    \underline u_{A,j}q_j
    &\le x_{A,j}\le\overline u_{A,j}q_j,
    \\
    \underline u_{B,j}q_j
    &\le x_{B,j}\le\overline u_{B,j}q_j.
    \label{eq:stage4-cell-perspective}
\end{align}
Using the Stage-3 first-order lifted kernel and certified second-order
remainder, define
\begin{align}
    \widetilde\Lambda_{st}^{\rm good}
    :=
    \sum_j
    K_\tau^{st}(c_j)
    \Big[
      q_j
      &+
      b_A^{st}(x_{A,j}-c_{A,j}q_j)
      \nonumber\\
      &+
      b_B^{st}(x_{B,j}-c_{B,j}q_j)
    \Big].
    \label{eq:stage4-lifted-overlap}
\end{align}
Let
\begin{equation}
    \rho_{st,j}
    :=
    \sup_{u\in C_j}
    \left|
      b_A^{st}(u_A-c_{A,j})
      +
      b_B^{st}(u_B-c_{B,j})
    \right|
\end{equation}
and
\begin{equation}
    \eta_{st,j}
    :=
    \min\!\left\{
      \frac{\rho_{st,j}^2}{2}e^{\rho_{st,j}},
      2+\rho_{st,j}
    \right\}.
\end{equation}
Then every continuum-feasible point satisfies
\begin{equation}
    \boxed{
    \left|
      \sum_zG[(z,s),(z,t)]
      -
      \widetilde\Lambda_{st}^{\rm good}
    \right|
    \le
    \sum_jq_j\eta_{st,j}
    +
    \beta.
    }
    \label{eq:stage4-certified-soc}
\end{equation}
Together with
\[
    G\succeq0,
    \qquad
    \sum_jq_j+\beta=1,
    \qquad
    0\le\beta\le\epsilon_{\rm mon},
\]
the perspective constraints, observed diagonal probabilities, and the A7
phase-error objective, this is the certified finite Stage-4 SDP implemented in
\texttt{phase/pnp-stage-4.py}.

The feasible-set direction is
\begin{equation}
    \boxed{
    \mathcal F_{\rm physical\ satisfying\ monitor\ interface}
    \subseteq
    \mathcal F_{\rm Stage\,4\ continuum}
    \subseteq
    \mathcal F_{\rm Stage\,4\ lifted}.
    }
    \label{eq:stage4-feasible-chain}
\end{equation}

\subsection{Observed statistics, honest simulation, and rate normalization}

For a real security calculation,
\begin{equation}
    P_{\rm pass|mon}
    :=
    \Pr(\text{central pass}\mid A_{\rm mon})
    \label{eq:stage4-ppass-observed}
\end{equation}
is an observed protocol statistic and should be supplied to the SDP, together
with the corresponding accepted-ensemble bit-error statistics.  It should
\emph{not} be generated from the worst-case energy distribution found by the
SDP.

The frequently used expression
\begin{equation}
    P_{\rm pass|mon}^{\rm hon}
    \simeq
    1-e^{-\eta\bar n_{\rm hon}}+p_{\rm dc}
    \label{eq:stage4-honest-pass}
\end{equation}
belongs only to a synthetic honest-performance model.  Here
\(\bar n_{\rm hon}\) denotes the honest outgoing brightness.  In particular,
the fact that a worst-case Stage-4 optimizer may push
\(\sum_jx_{A,j}\) or \(\sum_jx_{B,j}\) close to the cap does not mean that the
honest source has that brightness.

The key rate per monitor-accepted round is
\begin{equation}
    R_{\rm acc}
    =
    P_{\rm pass|mon}
    \left[
      1-h_2(e_{\rm ph})
      -f_{\rm ec}h_2(e_{\rm bit})
    \right].
    \label{eq:stage4-rate-accepted}
\end{equation}
The rate per original incoming round is
\begin{equation}
    \boxed{
    R_{\rm original}
    =
    P_{\rm mon}R_{\rm acc}.
    }
    \label{eq:stage4-rate-original}
\end{equation}
Equivalently, if the experimentally reported usable-event probability is
already
\[
    \Pr(A_{\rm mon}\wedge\text{central pass}),
\]
then \(P_{\rm mon}\) must not be multiplied a second time.

\subsection{Current numerical status and interpretation}

The current default exploratory parameters are
\[
    M=4,
    \qquad
    \tau=10^{-8},
    \qquad
    \bar n_{\rm hon}=10^{-2},
    \qquad
    u_{A,+}=u_{B,+}=1.2\,\frac{\bar n_{\rm hon}}{2},
    \qquad
    \epsilon_{\rm mon}=10^{-6}.
\]
With \(N_U=60\) cells per axis, the representative scan is
\begin{center}
\begin{tabular}{ccccc}
\toprule
\(L\) (km) & \(e_{\rm ph}\) & \(P_{\rm pass|mon}\) & \(R_{\rm acc}\) & \(\beta\)\\
\midrule
0  & 0.0393 & \(8.464\times10^{-3}\) & \(5.347\times10^{-3}\) & \(10^{-6}\)\\
2  & 0.0420 & \(7.722\times10^{-3}\) & \(4.782\times10^{-3}\) & \(10^{-6}\)\\
5  & 0.0462 & \(6.729\times10^{-3}\) & \(4.044\times10^{-3}\) & \(10^{-6}\)\\
10 & 0.0534 & \(5.349\times10^{-3}\) & \(3.049\times10^{-3}\) & \(10^{-6}\)\\
20 & 0.0706 & \(3.378\times10^{-3}\) & \(1.698\times10^{-3}\) & \(10^{-6}\)\\
40 & 0.1254 & \(1.346\times10^{-3}\) & \(4.389\times10^{-4}\) & \(10^{-6}\)\\
\bottomrule
\end{tabular}
\end{center}
These are conditional rates per monitor-accepted round.  They are not yet
end-to-end physical rates because \(P_{\rm mon}\) has not been supplied by a
physical monitor model.

Two optimizer features are informative.  First, \(\beta\) saturates
\(\epsilon_{\rm mon}\), showing that monitor leakage is directly
security-relevant.  Second, the good-component first moments tend to saturate
the allowed upper cap; without a separately certified conditional mean bound,
the worst-case SDP correctly lets Eve place accepted good mass as high as the
support constraint permits.

The present \(u_+\) scan should therefore be read as a
\emph{security-versus-cap} scan.  It does not yet optimize a physical detector
threshold, because changing \(u_+\) in the current code does not simultaneously
recompute \(P_{\rm mon}\), the honest accepted brightness, or
\(P_{\rm pass|mon}\) from a detector/source model.

\subsection{Final rigour statement for Stage 4}

Stage~4 is rigorous conditional on
Eq.~\eqref{eq:stage4-monitor-interface}.  It establishes that a sufficiently
tight accepted-round upper-tail constraint removes the specific rare-bright
concentration freedom that defeats the mean-only analysis.

Stage~4 does \emph{not} establish that a particular beam splitter,
photodiode, PNA, or comparator circuit satisfies that interface.  It also does
not require a QND measurement, and it should not be read as a physical
implementation model.  Deriving the interface from a concrete destructive
monitor is the purpose of Stage~5.

\section{Stage 5: physical PNA realisation}
\label{sec:stage5-plan}

Stage~5 will replace the abstract Stage-4 monitor interface by an explicit
trusted photon-number-analyzer model.  Existing plug-and-play source-monitor
work, in particular the APN/PNA/PND framework of Peng, Xu, and Guo
(arXiv:0908.1641), shows that the relevant monitoring hardware can be built
from ordinary destructive photodetection on a beam-splitter arm; no QND
measurement is required.  Our security use of the monitor will differ in that
the PNA outcome is to be incorporated directly into the accepted-round
untrusted-source Gram analysis.

\subsection{Stage 5A: joint PNA--beam-splitter model}
\label{sec:stage5a}

Stage~5A treats the trusted beam splitter and the destructive PNA as one
quantum instrument.  There is no separate attenuation element in this stage.
The beam splitter itself both sends a large fraction of the incoming light to
the monitor and leaves the complementary output as the signal mode.  The
secret phase encoding is applied only after the monitor decision.

For one arm define
\[
    M=\text{incoming photon number},\qquad
    C=\text{monitor-port photon number},\qquad
    N=\text{surviving signal photon number}.
\]
Let \(\tau_{\rm s}\) be the signal-port transmissivity, so that the monitor
fraction is \(1-\tau_{\rm s}\).  For an ideal lossless beam splitter,
\(M=C+N\).

\subsubsection{Exact beam-splitter map}

For an input Fock state \(|m\rangle\), with vacuum in the unused port,
\begin{equation}
    |m\rangle
    \longmapsto
    \sum_{n=0}^{m}
    \sqrt{\binom mn}
    \tau_{\rm s}^{n/2}
    (1-\tau_{\rm s})^{(m-n)/2}
    |n\rangle_S|m-n\rangle_D .
    \label{eq:stage5a-bs-map}
\end{equation}
The incoming state may be an arbitrary
\(\rho_{ME_0}\), including Fock coherences and entanglement with Eve.

\subsubsection{General phase-insensitive PNA acceptance}

A phase-insensitive acceptance rule is described by the diagonal POVM effect
\begin{equation}
    E_{\rm acc}
    =
    \sum_{c=0}^{\infty}a_c|c\rangle\!\langle c|,
    \qquad
    0\le a_c\le1,
    \label{eq:stage5a-accept-effect}
\end{equation}
where
\[
    a_c
    :=
    \Pr(A_{\rm mon}\mid C=c).
\]
The response coefficients \(a_c\) are where a later physical detector model
will enter: detector efficiency, finite resolution, electronic noise,
thresholds, saturation, and related trusted imperfections all modify this
response.

The unnormalised accepted signal--Eve state is
\begin{equation}
\begin{aligned}
    \widetilde\sigma_{\rm acc}^{SE_0}
    :=
    \operatorname{Tr}_{D}
    \Big[
      (I_{SE_0}\otimes E_{\rm acc})
      B_{\tau_{\rm s}}
      \big(
        \rho_{ME_0}\otimes|0\rangle\!\langle0|
      \big)
      B_{\tau_{\rm s}}^\dagger
    \Big].
    \label{eq:stage5a-accepted-state}
\end{aligned}
\end{equation}
Define
\begin{equation}
    P_{\rm mon}
    :=
    \operatorname{Tr}\widetilde\sigma_{\rm acc}^{SE_0},
    \qquad
    \sigma_{\rm acc}^{SE_0}
    :=
    \frac{\widetilde\sigma_{\rm acc}^{SE_0}}{P_{\rm mon}}.
    \label{eq:stage5a-normalized-state}
\end{equation}
Equation~\eqref{eq:stage5a-accepted-state} is the fundamental Stage-5A map.
The attenuation produced by the beam splitter and the PNA outcome are
correlated and must be treated jointly.

\subsubsection{Accepted surviving photon-number distribution}

Although the incoming state can contain arbitrary coherences, the accepted
signal photon-number diagonal depends only on
\begin{equation}
    p_m
    :=
    \operatorname{Tr}_{E_0}
    \langle m|\rho_{ME_0}|m\rangle.
    \label{eq:stage5a-input-diagonal}
\end{equation}
The unnormalised probability for an accepted signal to contain \(n\) photons is
\begin{equation}
    \boxed{
    \widetilde q_n
    =
    \sum_{m=n}^{\infty}
      p_m
      \binom mn
      \tau_{\rm s}^{n}
      (1-\tau_{\rm s})^{m-n}
      a_{m-n}.
    }
    \label{eq:stage5a-qtilde}
\end{equation}
Equivalently, setting \(c=m-n\),
\begin{equation}
    \widetilde q_n
    =
    \sum_{c=0}^{\infty}
      p_{n+c}
      \binom{n+c}{n}
      \tau_{\rm s}^{n}
      (1-\tau_{\rm s})^{c}
      a_c.
    \label{eq:stage5a-qtilde-c}
\end{equation}
Hence
\begin{equation}
    P_{\rm mon}
    =
    \sum_n\widetilde q_n,
    \qquad
    q_n^{\rm acc}
    =
    \frac{\widetilde q_n}{P_{\rm mon}}.
    \label{eq:stage5a-qacc}
\end{equation}
Thus \(q_n^{\rm acc}\) is the actual photon-number distribution in the
surviving signal mode conditioned on monitor acceptance.

\subsubsection{Common phase orbit and exact conditional overlap}

The secret phase encoding is applied only after the monitor decision:
\[
    U_\phi=e^{i\phi\hat n_S}.
\]
Therefore
\begin{equation}
    \boxed{
    \rho_\phi^{\rm acc}
    =
    U_\phi\sigma_{\rm acc}U_\phi^\dagger.
    }
    \label{eq:stage5a-common-orbit}
\end{equation}
The common phase orbit is automatic; no additional attenuator is needed to
create or restore it.

For two phase labels \(s,t\),
\begin{align}
    \Lambda_{st}^{\rm acc}
    &=
    \operatorname{Tr}
    \left[
      \sigma_{\rm acc}
      e^{i\Delta\phi^{st}\hat n_S}
    \right]
    \nonumber\\
    &=
    \boxed{
    \sum_{n=0}^{\infty}
      q_n^{\rm acc}
      e^{i\Delta\phi^{st}n}.
    }
    \label{eq:stage5a-overlap-q}
\end{align}
As in Stage~1, Fock coherences drop out because the phase-difference operator
is diagonal in photon number.

The same overlap can be written directly in terms of the incoming
distribution.  Define
\begin{equation}
    F_m(\theta)
    :=
    \sum_{c=0}^{m}
      a_c
      \binom mc
      (1-\tau_{\rm s})^c
      \tau_{\rm s}^{m-c}
      e^{i\theta(m-c)}
    \label{eq:stage5a-Fm}
\end{equation}
and
\begin{equation}
    A_m
    :=
    F_m(0)
    =
    \Pr(A_{\rm mon}\mid M=m).
    \label{eq:stage5a-Am}
\end{equation}
Then
\begin{equation}
    \boxed{
    \Lambda^{\rm acc}(\theta)
    =
    \frac{
      \sum_m p_mF_m(\theta)
    }{
      \sum_m p_mA_m
    }.
    }
    \label{eq:stage5a-direct-overlap}
\end{equation}

For Alice and Bob, let the arbitrary joint input diagonal be
\(p_{m_A,m_B}\).  No product-source assumption is made.  For independent
trusted local beam splitters and monitor responses,
\begin{equation}
    P_{\rm mon}
    =
    \sum_{m_A,m_B}
      p_{m_A,m_B}
      A^A_{m_A}A^B_{m_B},
    \label{eq:stage5a-pmon-twoarm}
\end{equation}
and
\begin{equation}
    \boxed{
    \Lambda_{st}^{\rm acc}
    =
    \frac{
      \sum_{m_A,m_B}
      p_{m_A,m_B}
      F^A_{m_A}(\Delta\phi_A^{st})
      F^B_{m_B}(\Delta\phi_B^{st})
    }{
      P_{\rm mon}
    }.
    }
    \label{eq:stage5a-overlap-twoarm}
\end{equation}

\subsubsection{Direct tail certification on the surviving signal}

Stage~5A can formulate the security cutoff directly in terms of the
\emph{actual} surviving signal photon number \(N\).  For a chosen cutoff
\(N_+\), define
\begin{equation}
\begin{aligned}
    H_m(N_+)
    &:=
    \Pr(A_{\rm mon},\,N>N_+\mid M=m)
    \\
    &=
    \sum_{n=N_++1}^{m}
      \binom mn
      \tau_{\rm s}^{n}
      (1-\tau_{\rm s})^{m-n}
      a_{m-n}.
    \label{eq:stage5a-Hm}
\end{aligned}
\end{equation}
The source-independent worst case is
\begin{equation}
    \boxed{
    \delta(N_+)
    :=
    \sup_{m\ge0}H_m(N_+).
    }
    \label{eq:stage5a-delta}
\end{equation}
For any Eve-supplied input distribution,
\begin{equation}
    \Pr(A_{\rm mon},N>N_+)
    \le
    \delta(N_+),
    \label{eq:stage5a-unconditional-tail}
\end{equation}
and hence, using the observed monitor acceptance probability,
\begin{equation}
    \boxed{
    \Pr(N>N_+\mid A_{\rm mon})
    \le
    \frac{\delta(N_+)}{P_{\rm mon}}.
    }
    \label{eq:stage5a-conditional-tail}
\end{equation}
For two arms a simple conservative union bound gives
\begin{equation}
    \Pr(
      N_A>N_{A,+}\ \text{or}\ N_B>N_{B,+}
      \mid A_{\rm mon}
    )
    \le
    \frac{
      \delta_A(N_{A,+})+\delta_B(N_{B,+})
    }{
      P_{\rm mon}
    }.
    \label{eq:stage5a-twoarm-tail}
\end{equation}
A tighter joint bound can be studied later if necessary.

\subsubsection{Ideal upper-threshold monitor: first architecture test}

Before including detector inefficiency and noise, the first sanity check is an
ideal photon-number-resolving upper-threshold monitor,
\begin{equation}
    \boxed{
    a_c
    =
    \mathbf 1[c\le c_{\max}].
    }
    \label{eq:stage5a-ideal-response}
\end{equation}
This is an ideal \emph{upper-threshold} photon-number monitor rather than the
full two-threshold PNA of Peng, Xu, and Guo.  Only the upper side is needed to
test the concentration mechanism.

For this ideal response,
\begin{equation}
    H_m(N_+)
    =
    \Pr(
      C\le c_{\max},\,N>N_+
      \mid M=m
    ).
    \label{eq:stage5a-ideal-H}
\end{equation}
Since \(C=m-N\), this can be written as
\begin{equation}
    H_m(N_+)
    =
    \sum_{n=\max\{N_++1,m-c_{\max}\}}^{m}
    \binom mn
    \tau_{\rm s}^{n}
    (1-\tau_{\rm s})^{m-n}.
    \label{eq:stage5a-ideal-H-sum}
\end{equation}

In this ideal case the supremum over the arbitrary incoming Fock number can
actually be located exactly.  Let
\begin{equation}
    m_*:=c_{\max}+N_++1.
    \label{eq:stage5a-mstar}
\end{equation}
For \(m\le m_*\), the monitor condition is no more restrictive than the
condition \(N>N_+\), so
\[
    H_m(N_+)=\Pr[\operatorname{Bin}(m,\tau_{\rm s})>N_+],
\]
which is nondecreasing in \(m\).  For \(m\ge m_*\), the event
\(C\le c_{\max}\) already implies \(N>N_+\), so
\[
    H_m(N_+)=\Pr[\operatorname{Bin}(m,1-\tau_{\rm s})\le c_{\max}],
\]
which is nonincreasing in \(m\).  Therefore the maximum occurs at \(m_*\):
\begin{equation}
    \boxed{
    \delta_{\rm ideal}(N_+)
    =
    H_{m_*}(N_+)
    =
    \Pr[
      \operatorname{Bin}(m_*,\tau_{\rm s})>N_+
    ].
    }
    \label{eq:stage5a-ideal-delta}
\end{equation}
Equivalently,
\[
    \delta_{\rm ideal}(N_+)
    =
    \Pr[
      \operatorname{Bin}(m_*,1-\tau_{\rm s})
      \le c_{\max}
    ].
\]

This closed form is the first quantitative feasibility test for the one-BS
architecture.  For a very small signal fraction \(\tau_{\rm s}\), a monitor
threshold \(c_{\max}\) of order the bright incoming pulse can coexist with a
very small probability that more than \(N_+\) photons escape into the signal
port.  The point of the calculation is to verify this in a worst-case,
source-independent way rather than infer it only from the conditional mean
\(\tau_{\rm s}m\).

\subsubsection{Finite-Fock Gram reduction when the surviving cutoff is small}

If Eq.~\eqref{eq:stage5a-conditional-tail} gives a useful bound for modest
\(N_+\), no Stage-3 continuum variable is needed in Stage~5A.  Decompose the
accepted surviving distribution as
\[
    q^{\rm acc}
    =
    q^{\rm good}+q^{\rm bad},
    \qquad
    \beta\le\epsilon_5,
\]
with the good part supported on
\[
    0\le n_A\le N_{A,+},
    \qquad
    0\le n_B\le N_{B,+}.
\]
Then
\begin{equation}
\begin{aligned}
    \Lambda_{st}^{\rm acc}
    =
    \sum_{\substack{
      0\le n_A\le N_{A,+}\\
      0\le n_B\le N_{B,+}
    }}
      q_{n_A,n_B}^{\rm good}
      e^{
        i\Delta\phi_A^{st}n_A
        +i\Delta\phi_B^{st}n_B
      }
    +
    r_{st},
    \qquad
    |r_{st}|\le\beta.
    \label{eq:stage5a-finite-overlap}
\end{aligned}
\end{equation}
This can be inserted directly into the same Gram phase-error SDP.  There are
no cell first moments or Taylor remainders in this finite-Fock version.

The preliminary script \texttt{phase/pnp-stage-5a.py} is intended first to
implement the exact ideal-threshold quantity
Eq.~\eqref{eq:stage5a-ideal-delta}, together with the general beam-splitter/PNA
response functions above.  It should be treated as exploratory until a
realistic detector response \(a_c\) and a corresponding certified
worst-case tail calculation are supplied.

\subsection{Stage 5B: PNA--beam splitter plus additional trusted attenuation}

Stage~5B will append an optional trusted pure-loss channel of transmissivity
\(\lambda\) to the Stage-5A surviving signal before phase encoding.  This is
the more general architecture and may be experimentally convenient because
the monitor splitting ratio and final QKD brightness can then be tuned
independently.

The additional attenuator is \emph{not} required to create or restore the
common phase orbit.  Its role is only to provide extra controllable signal
loss and, analytically, an additional known thinning map.  Setting
\(\lambda=1\) recovers Stage~5A exactly.

\section{Open points}

The main unresolved items are now:

\begin{enumerate}
    \item \textbf{Stage 5A: derive the joint PNA--BS instrument.}
    Starting from an arbitrary Eve-supplied state, model the trusted
    beam splitter, destructive monitor POVM, detector efficiency, finite
    resolution, additive noise, saturation, timing and mode dependence.
    Derive a certified conditional constraint on the surviving signal state
    and connect it directly to the Gram overlaps.

    \item \textbf{Stage 5B: add optional trusted attenuation.}
    Extend Stage~5A by a subsequent pure-loss channel of transmissivity
    \(\lambda\).  Verify explicitly that \(\lambda=1\) reduces to Stage~5A and
    determine whether the extra attenuation materially improves the
    security/performance tradeoff or merely implementation flexibility.

    \item \textbf{Physical threshold/performance optimisation.}
    Once the monitor model is available, jointly study monitor splitting
    ratio, PNA thresholds, false-accept leakage, \(P_{\rm mon}\), honest
    accepted brightness, and \(P_{\rm pass|mon}\).  This is the calculation
    needed to turn the current Stage-4 security-versus-cap scan into an
    end-to-end key-rate optimisation.

    \item \textbf{Stage-4 numerical convergence and solver accuracy.}
    Quantify cell-refinement error and rerun publication-quality Stage-4
    scans with MOSEK or another sufficiently accurate solver.  The current
    \(\epsilon_{\rm mon}=10^{-6}\) value must continue to be labelled as a
    placeholder until Stage~5 supplies a detector-derived value.

    \item \textbf{Revisit the number of test bases.}
    The apparent saturation near \(M=4\) was originally found in a regime
    dominated by the concentration attack.  Once the physical PNA model is in
    place, the optimal number of phase bases should be re-optimised.
\end{enumerate}

\section{Working conclusion}

The present results separate the source-monitoring problem into three levels.
A first-moment/APN-type guarantee is insufficient: rare bright pulses remain
compatible with the monitored mean, and channel loss lets those rare
components dominate the accepted central detections.  Strong trusted
attenuation changes the overlap kernel and enables the continuum formulation,
but does not by itself eliminate that concentration freedom.

An accepted-round upper-tail/PNA-type guarantee is sufficient at the abstract
level.  Stage~4 proves this statement conditionally through the Gram SDP:
once all but an \(\epsilon_{\rm mon}\) fraction of accepted effective-energy
mass is confined to a fixed compact region, the rare-bright mechanism is
removed and useful phase-error bounds are recovered in the present scans.

What remains is the physical reduction.  Stage~5A will derive the accepted
signal state produced by a joint trusted beam splitter and destructive PNA,
with the same beam splitter allowed to provide the signal attenuation.
Stage~5B will add optional further trusted attenuation.  The common phase
orbit is not the obstacle in either case, because monitoring and attenuation
occur before the secret phase encoding; the real task is to derive the
accepted-state photon-number constraints and overlap kernel from the physical
monitor instrument.

\appendix

\section{Clarifications on the trusted-attenuation construction}
\label{app:trusted-attenuation-clarifications}

This appendix collects three points that are easy to conflate in the main
derivation:

\begin{enumerate}
    \item mixed states are not by themselves an obstruction to a Gram
    construction; the useful extra structure is a known common unitary orbit;
    \item the auxiliary purifier used in the Gram relaxation is not the
    physical reflected loss port, and is not physically assumed to be in
    Eve's possession;
    \item the mean-energy constraint must be placed on the variable
    corresponding to the location of the trusted monitor.  When the monitored
    quantity is the light leaving Alice's and Bob's laboratories after trusted
    attenuation, the same bound appears with a factor \(1/\tau\) when written
    in terms of the pre-attenuation photon-number distribution.
\end{enumerate}

\subsection{Common orbit versus a generic mixed-state family}
\label{app:common-orbit-example}

The statement that loss makes the states mixed and therefore prevents any
Gram construction is too coarse.  The relevant distinction is instead
\[
    \boxed{
    \text{mixed states on a known common unitary orbit}
    \quad\text{versus}\quad
    \text{a generic mixed-state family}.
    }
\]
A simple three-state example makes this explicit.

\subsubsection{Three phase states that remain on one orbit}

Work in the single-rail subspace
\(\operatorname{span}\{|0\rangle,|1\rangle\}\), and choose the three phases
\[
    \phi_j=\frac{2\pi j}{3},
    \qquad
    j=0,1,2.
\]
Define
\[
    U_j
    =
    |0\rangle\!\langle0|
    +
    e^{i\phi_j}|1\rangle\!\langle1|,
\]
and start from
\[
    |+\rangle
    =
    \frac{|0\rangle+|1\rangle}{\sqrt2}.
\]
Before loss,
\[
    |\psi_j\rangle
    =
    U_j|+\rangle
    =
    \frac{|0\rangle+e^{i\phi_j}|1\rangle}{\sqrt2}.
\]

Let \(\mathcal L_\eta\) be pure loss of transmissivity \(\eta\).  Since pure
loss is phase covariant,
\[
    \mathcal L_\eta(U_j\rho U_j^\dagger)
    =
    U_j\mathcal L_\eta(\rho)U_j^\dagger.
\]
Therefore the three output states satisfy
\[
    \rho_j
    =
    U_j\rho_0U_j^\dagger.
\]
Explicitly,
\begin{equation}
    \rho_j
    =
    \begin{pmatrix}
        1-\eta/2
        &
        \frac{\sqrt{\eta}}{2}e^{-i\phi_j}
        \\[1mm]
        \frac{\sqrt{\eta}}{2}e^{i\phi_j}
        &
        \eta/2
    \end{pmatrix}.
    \label{eq:app-phase-orbit-rho}
\end{equation}
All three states have the same spectrum, as they must under unitary
conjugation.

Choose one purification
\[
    |\chi_0\rangle_{TR}
\]
of \(\rho_0\), and define
\[
    |\chi_j\rangle
    =
    (U_j\otimes I_R)|\chi_0\rangle.
\]
Then
\[
    \operatorname{Tr}_R|\chi_j\rangle\!\langle\chi_j|
    =
    \rho_j,
\]
and the pairwise overlaps are fixed by
\begin{align}
    \langle\chi_j|\chi_k\rangle
    &=
    \operatorname{Tr}
    \left[
      \rho_0U_j^\dagger U_k
    \right]
    \nonumber\\
    &=
    1-\frac{\eta}{2}
    +
    \frac{\eta}{2}
    e^{i(\phi_k-\phi_j)}.
    \label{eq:app-phase-orbit-overlap}
\end{align}
Thus the three mixed output states admit one consistent purified Gram matrix.

This is exactly the structure used in our trusted-attenuation argument.  More
generally, if a lossy mode-discrimination family happens to satisfy
\[
    \rho_j=W_j\rho_0W_j^\dagger
\]
for known \(W_j\), then the same common-purification construction can be made.
The observation in arXiv:2012.11104 that their generic lossy
channel-discrimination problem is no longer directly described by their
original pure-state SDP should therefore not be read as a no-go theorem for
every special covariant family.

For a discrimination problem, however, the purification-based Gram
optimization need not equal the physical mixed-state discrimination optimum:
the physical receiver acts only on \(T\), whereas an unconstrained Gram
optimization can effectively exploit additional freedom associated with the
purifier.  It is therefore naturally interpreted as an upper bound unless the
restriction to physical measurements on \(T\) is enforced separately.

\subsubsection{Three states that do not remain on one orbit}

Now take instead
\[
    |\psi_0\rangle=|0\rangle,
    \qquad
    |\psi_1\rangle=|1\rangle,
    \qquad
    |\psi_2\rangle
    =
    \frac{|0\rangle+|1\rangle}{\sqrt2}.
\]
Before loss these three pure states can certainly be related to
\(|0\rangle\) by suitable unitaries.  After the same pure-loss channel,
however,
\begin{align}
    \rho_0
    &=
    \begin{pmatrix}
        1&0\\
        0&0
    \end{pmatrix},
    \\
    \rho_1
    &=
    \begin{pmatrix}
        1-\eta&0\\
        0&\eta
    \end{pmatrix},
    \\
    \rho_2
    &=
    \begin{pmatrix}
        1-\eta/2&\sqrt{\eta}/2\\
        \sqrt{\eta}/2&\eta/2
    \end{pmatrix}.
    \label{eq:app-nonorbit-rhos}
\end{align}
For generic \(0<\eta<1\), these states have different spectra.  Hence there
cannot exist unitaries \(W_j\) such that
\[
    \rho_j=W_j\rho_0W_j^\dagger,
\]
because unitary conjugation preserves eigenvalues.

One can of course purify each \(\rho_j\) separately, but then there is no
distinguished relation of the form
\[
    |\chi_j\rangle
    =
    (W_j\otimes I_R)|\chi_0\rangle
\]
that fixes all pairwise overlaps from one reference state.  Different choices
of purification can give different overlaps.  This is the obstruction that
the common-orbit property removes.

\subsection{What the auxiliary purification means physically}
\label{app:purifier-meaning}

It is important to distinguish four systems:
\[
    T,\qquad E_0,\qquad L,\qquad R.
\]
Here \(T\) is the transmitted light returned to the untrusted node,
\(E_0\) is Eve's ancilla from preparing the source, \(L\) is the physical
trusted loss port of the attenuator, and \(R\) is an auxiliary mathematical
purifier introduced only after \(L\) has been traced out.

The physical post-attenuation state is
\[
    \rho_s^{TE_0}
    =
    \operatorname{Tr}_{L}
    |\Phi_s\rangle\!\langle\Phi_s|_{TE_0L}.
\]
The system \(L\) is trusted and unavailable to Eve.  By phase covariance,
\[
    \rho_s^{TE_0}
    =
    W_s\rho_0^{TE_0}W_s^\dagger.
\]
We then choose a mathematical purification
\[
    |\chi_0\rangle_{TE_0R}
\]
of \(\rho_0^{TE_0}\) and define
\[
    |\chi_s\rangle
    =
    (W_s\otimes I_R)|\chi_0\rangle.
\]

A physical Eve acts only on \(TE_0\).  If her instrument has Kraus operators
\(K_z\), then on the purified representation the physical residual vectors
are
\begin{equation}
    |e_s^z\rangle
    =
    (K_z\otimes I_R)|\chi_s\rangle.
    \label{eq:app-physical-purified-attack}
\end{equation}
Thus the introduction of \(R\) does \emph{not} physically assume that Eve
possesses the purifier.

The outer relaxation occurs only afterwards.  The SDP keeps the Gram
positivity, observed diagonal entries, and overlap-conservation constraints,
but does not explicitly enforce that all residual vectors arise from one
common instrument of the form \(K_z\otimes I_R\).  Consequently,
\[
    \mathcal F_{\rm physical}
    \subseteq
    \mathcal F_{\rm Gram}.
\]
One may heuristically interpret the enlarged Gram problem as giving Eve access
to additional purifier-like degrees of freedom, but this is a mathematical
upper-bounding device rather than a statement about the physical apparatus.

\subsubsection{Why \(R\) must not be identified with the physical loss port}

If the actual reflected port \(L\) were retained, the complete beamsplitter
evolution would be unitary and independent of the encoded label \(s\).  Hence
\begin{equation}
    \langle\Phi_s|\Phi_t\rangle
    =
    \langle\psi_s|\psi_t\rangle.
    \label{eq:app-physical-loss-port-overlap}
\end{equation}
At this global level the attenuator does not change distinguishability:
nothing has actually been discarded.

By contrast, after tracing out \(L\) and constructing the covariant
purification of the reduced transmitted state, the overlap becomes
\begin{equation}
    \langle\chi_s|\chi_t\rangle
    =
    \operatorname{Tr}
    \left[
      \rho_0^{TE_0}W_s^\dagger W_t
    \right],
    \label{eq:app-virtual-purifier-overlap}
\end{equation}
which yields the attenuated kernel
\[
    \sum_m p_m
    \left(
      1-\tau+\tau e^{i\Delta\phi}
    \right)^m
\]
in one mode.

For example, for
\[
    |\psi_\phi\rangle
    =
    \frac{|0\rangle+e^{i\phi}|1\rangle}{\sqrt2},
\]
keeping the physical beamsplitter environment gives
\[
    \langle\Phi_\phi|\Phi_{\phi'}\rangle
    =
    \frac12
    \left(
      1+e^{i(\phi'-\phi)}
    \right),
\]
independent of \(\tau\).  The common purification of the reduced transmitted
state instead gives
\[
    \langle\chi_\phi|\chi_{\phi'}\rangle
    =
    1-\frac{\tau}{2}
    +
    \frac{\tau}{2}e^{i(\phi'-\phi)}.
\]
For opposite phases,
\[
    \langle\chi_\phi|\chi_{\phi+\pi}\rangle
    =
    1-\tau,
\]
whereas the global state including the physical loss port remains orthogonal.
This explicitly shows why
\[
    \boxed{R\neq L.}
\]

\subsection{Where the average-photon-number bound belongs}
\label{app:energy-bound-placement}

The placement of the energy constraint is fixed by the physical location of
the trusted monitor and by the variable used in the optimization.

Let
\[
    M_X
\]
denote the photon number entering the trusted attenuator in arm \(X\), and
let
\[
    N_X
\]
denote the photon number leaving it toward the untrusted node.  Pure loss of
transmissivity \(\tau\) satisfies exactly
\begin{equation}
    \mathbb E[N_X]
    =
    \tau\mathbb E[M_X].
    \label{eq:app-loss-first-moment}
\end{equation}

The convention used in the present notes is that
\(\bar n\) is the total monitored mean photon number \emph{leaving} Alice's
and Bob's laboratories after the trusted attenuation, with the symmetric
per-arm bounds
\begin{equation}
    \mathbb E[N_A]
    \le
    \frac{\bar n}{2},
    \qquad
    \mathbb E[N_B]
    \le
    \frac{\bar n}{2}.
    \label{eq:app-output-energy-bound}
\end{equation}
If Stage~2 is written in terms of the pre-attenuation distribution
\(p_{m_A,m_B}\), Eq.~\eqref{eq:app-loss-first-moment} makes these equivalent
to
\begin{equation}
    \sum_{m_A,m_B}
      m_Ap_{m_A,m_B}
    \le
    \frac{\bar n}{2\tau},
    \qquad
    \sum_{m_A,m_B}
      m_Bp_{m_A,m_B}
    \le
    \frac{\bar n}{2\tau}.
    \label{eq:app-input-energy-bound}
\end{equation}
Thus the factor \(1/\tau\) does not represent an additional assumption.  It
is simply the same post-attenuation monitored bound rewritten in the input
variables.

If instead the optimization is written in terms of the transmitted
photon-number distribution \(q_{n_A,n_B}\), the natural constraints are
directly
\[
    \sum_{n_A,n_B}n_Aq_{n_A,n_B}
    \le\frac{\bar n}{2},
    \qquad
    \sum_{n_A,n_B}n_Bq_{n_A,n_B}
    \le\frac{\bar n}{2}.
\]
Likewise, in Stage~3 the rescaled variables
\[
    u_X=\tau m_X
\]
give
\begin{equation}
    \int u_A\,d\nu
    \le\frac{\bar n}{2},
    \qquad
    \int u_B\,d\nu
    \le\frac{\bar n}{2}.
    \label{eq:app-continuum-energy-bound}
\end{equation}

Hence the same physical constraint appears in three equivalent coordinate
systems:
\begin{equation}
\boxed{
    \mathbb E[N_X]\le\frac{\bar n}{2}
    \quad\Longleftrightarrow\quad
    \mathbb E[M_X]\le\frac{\bar n}{2\tau}
    \quad\Longleftrightarrow\quad
    \mathbb E[U_X]\le\frac{\bar n}{2}.
}
    \label{eq:app-energy-coordinate-equivalence}
\end{equation}

If the actual trusted monitor were placed \emph{before} the attenuator instead,
then the certified physical statement would instead be a bound directly on
\(\mathbb E[M_X]\), and the factors of \(\tau\) would have to be moved
accordingly.  The formulas above therefore depend on the protocol convention
that the monitored quantity is the outgoing energy after trusted attenuation.

\subsubsection{Where the truncation tail charge belongs}

There is a second, related placement issue.  In Stage~2 the rigorous
photon-number cutoff is imposed on the \emph{input} variables:
\[
    S_N
    =
    \{(m_A,m_B):m_A+m_B\le N\}.
\]
If
\[
    \varepsilon_N
    =
    1-\sum_{S_N}p_{m_A,m_B},
\]
then every point in the discarded set obeys
\[
    m_A+m_B\ge N+1.
\]
Therefore the total input-energy relaxation
\begin{equation}
    \sum_{S_N}
    (m_A+m_B)p_{m_A,m_B}
    +(N+1)\varepsilon_N
    \le
    \frac{\bar n}{\tau}
    \label{eq:app-input-tail-energy}
\end{equation}
is sound.

The analogous statement is \emph{not} available for a tail defined only by
omitted output photon numbers.  A very large input Fock component can be
mapped by loss to any lower output photon number, including vacuum.  Thus
input tail mass is not concentrated in a high-photon-number output tail.

For the same reason, the common tail charge
\((N+1)\varepsilon_N\) cannot be added separately to both per-arm constraints:
a discarded point such as \((0,N+1)\) has zero Alice-arm photon number.  The
tail charge is valid for the combined photon number because only
\(m_A+m_B\) has a uniform lower bound on the triangular discarded set.

The Stage-3 cell relaxation follows the same logic.  For upper bounds on a
mean energy, each cell mass must be charged using a lower bound on the energy
inside that cell.  If
\[
    \underline u_{X,j}
    =
    \inf_{u\in C_j}u_X,
\]
then
\[
    \sum_j\underline u_{X,j}q_j
    \le
    \frac{\bar n}{2}
\]
is an outer relaxation.  Charging each cell at its centre without a
compensating error could instead exclude a physical measure and would therefore
be unsafe.


\end{document}
